% ==============================================================================
% ch08.tex -- Chapter 8: Coarse moduli schemes, cusps, and compactification
%   Source: book pp. 224--270 (physical pp. 235--281)
% ==============================================================================
\chapter{Coarse Moduli Schemes, Cusps, and Compactification}
\label{ch:8}
\markboth{8.\ \MakeUppercase{Coarse Moduli Schemes, Cusps, and Compactification}}{}

% ------------------------------------------------------------------------------
% Section 8.1
% ------------------------------------------------------------------------------
\section{Coarse moduli schemes}
\label{sec:8.1}

\kmitem\label{item:8.1.1}
Let $R$ be a ring, $\mathcal{P}$ a relatively representable moduli problem on \EllR\ which is affine over \EllR. We will recall the definition of the ``coarse moduli scheme'' $\Mcoarse(\mathcal{P})$. This is an $R$-scheme which agrees with $\Mfrak(\mathcal{P})$ for representable $\mathcal{P}$, and which is a ``best replacement'' if $\mathcal{P}$ is not representable.

To define $\Mcoarse(\mathcal{P})$ as an $R$-scheme, it suffices to do so locally on $R$. So we may assume that some integer $N \ge 3$ is invertible in $R$. Let $\mathcal{S}$ be any representable moduli problem which is finite etale galois over \EllR, with galois group $G$ (e.g., $\mathcal{S} = [\Gamma(N)]$, $G = \GL(2, \Z/N\Z)$). We define $\Mcoarse(\mathcal{P})$ as the quotient scheme
\[
\Mcoarse(\mathcal{P}) = \Mfrak(\mathcal{P}, \mathcal{S})/G .
\]
The resulting $R$-scheme is clearly independent of the auxiliary choice of $\mathcal{S}$ (and so patches together). It ``exists'' because $\Mfrak(\mathcal{P}, \mathcal{S})$ is itself affine.

\begin{klem}\label{lem:8.1.2}
If $\mathcal{P}$ is normal, then $\Mcoarse(\mathcal{P})$ is normal.
\end{klem}

\begin{proof}
If $\mathcal{P}$ is normal, then $\Mfrak(\mathcal{P}, \mathcal{S})$ is normal, hence also its quotient by a finite group.
\end{proof}

\kmitem\label{item:8.1.3}
Let $S$ be an $R$-scheme, $E/S$ an elliptic curve, and $a \in \mathcal{P}(E/S)$ a level $\mathcal{P}$-structure on $E/S$. Then there is a canonical ``classifying map'' of $R$-schemes
\[
S \to \Mcoarse(\mathcal{P})
\]
defined as follows. Locally on $S$, we may assume some integer $N \ge 3$ is invertible on $S$. Let $S_N$ denote the finite etale galois $S$-scheme
\[
S_N = [\Gamma(N)]_{E/S} ,
\]
$G$ the galois group $\GL(2, \Z/N\Z)$ of $S_N$ over $S$. Then $E \times_S S_N/S_N$ carries both a level $\mathcal{P}$-structure and a $[\Gamma(N)]$-structure, so we have an honest $G$-equivariant classifying map
\[
G\ \left(
\begin{tikzcd}[column sep=huge, row sep=large]
S_N \arrow[r] \arrow[d] & \Mfrak(\mathcal{P}, \Gamma(N)) \arrow[d] \\
S \arrow[r, dashed] & \Mcoarse(\mathcal{P}) ,
\end{tikzcd}
\right.
\]
which therefore induces a map between the quotients by the $G$-action. Again, one verifies easily that the map $S \to \Mcoarse(\mathcal{P})$ constructed this way is independent of the auxiliary choice of $N$. Of course if $\mathcal{P}$ is representable, this map $S \to \Mcoarse(\mathcal{P}) = \Mfrak(\mathcal{P})$ is the classifying map for the isomorphism class of the data $(E/S/R, a)$.

\begin{customlem}{8.1.3.1}\label{lem:8.1.3.1}
When $k$ is an algebraically closed field, then
\[
\begin{aligned}
\Mcoarse(\mathcal{P})(k) =\ &\text{the set of $k$-isomorphism classes of ``elliptic curves} \\
&\text{$E/k$ with given level $\mathcal{P}$-structure.''}
\end{aligned}
\]
\end{customlem}

\begin{proof}
By (\ref{cor:A7.2.2}), on $k$-valued points we have a bijection
\[
\Mfrak(\mathcal{P}, [\Gamma(N)])(k)/G \xrightarrow{\ \sim\ } \Mcoarse(\mathcal{P})(k) . \qedhere
\]
\end{proof}

\kmitem\label{item:8.1.4}
If $\mathcal{P}$ and $\mathcal{P}'$ are two relatively representable moduli problems on \EllR, any morphism
\[
\mathcal{P} \to \mathcal{P}'
\]
over \EllR\ induces a morphism of $R$-schemes
\[
\Mcoarse(\mathcal{P}) \to \Mcoarse(\mathcal{P}') ,
\]
which, for any representable $\mathcal{S}$ which is finite etale galois over \EllR\ with galois group $G$, sits in a commutative diagram
\[
\begin{tikzcd}[column sep=huge, row sep=large]
\Mfrak(\mathcal{P}, \mathcal{S}) \arrow[r] \arrow[d, "{\text{divide by }G}"] & \Mfrak(\mathcal{P}', \mathcal{S}) \arrow[d, "{\text{divide by }G}"] \\
\Mcoarse(\mathcal{P}) \arrow[r, dashed] & \Mcoarse(\mathcal{P}') .
\end{tikzcd}
\]

\begin{klem}\label{lem:8.1.5}
Suppose that a finite group $G$ acts on $\mathcal{P}$. Then
\[
\Mcoarse(\mathcal{P})/G \xrightarrow{\ \sim\ } \Mcoarse(\mathcal{P}/G) .
\]
\end{klem}

\begin{proof}
The action of $G$ on $\mathcal{P}$ induces an action of $G$ on $\Mcoarse(\mathcal{P})$. The morphism $\mathcal{P} \to \mathcal{P}/G$ induces a $G$-equivariant map
\[
\Mcoarse(\mathcal{P}) \to \Mcoarse(\mathcal{P}/G) ,
\]
whence a map
\[
\Mcoarse(\mathcal{P})/G \to \Mcoarse(\mathcal{P}/G) .
\]
This is the map we must show to be an isomorphism. The question is local on $R$, so we may assume some odd prime $\ell$ invertible in $R$. Then by definition we have
\begin{align*}
\Mcoarse(\mathcal{P}/G) &= \Mfrak(\mathcal{P}/G, [\Gamma(\ell)])/\GL(2, \F_\ell) \\
&= (\Mfrak(\mathcal{P}, [\Gamma(\ell)])/G)/\GL(2, \F_\ell) \\
&= \Mfrak(\mathcal{P}, [\Gamma(\ell)])/G \times \GL(2, \F_\ell) \\
&= (\Mfrak(\mathcal{P}, [\Gamma(\ell)])/\GL(2, \F_\ell))/G \\
&= \Mcoarse(\mathcal{P})/G . \qedhere
\end{align*}
\end{proof}

\begin{kprop}\label{prop:8.1.6}
For any extension of scalars $R \to R'$, we have a canonical morphism of $R'$-schemes
\[
\Mcoarse(\mathcal{P} \tens{R} R') \to \Mcoarse(\mathcal{P}) \tens{R} R' .
\]
This morphism is an isomorphism if any of the following conditions holds
\begin{enumerate}[(1)]
\item $\mathcal{P}$ is representable.
\item $R \to R'$ is flat.
\item the integer $6 = 2 \times 3$ is invertible in $R$.
\item $\mathcal{P}$ is the quotient of a representable problem $\mathcal{P}'$ by a finite group $G$ whose order is invertible in $R'$.
\end{enumerate}
\end{kprop}

\begin{proof}
If $\mathcal{P}$ is representable, then $\Mcoarse(\mathcal{P}) = \Mfrak(\mathcal{P})$ and the assertion is tautologous. If (4) holds, simply apply (\ref{prop:A7.1.3}, (3)) to the action of $G$ on the coordinate ring of $\Mfrak(\mathcal{P}')$. If $6$ is invertible, we can take for $\mathcal{P}'$ the simultaneous problem $(\mathcal{P}, [\Gamma(3)])$, with $G = \GL(2, \F_3)$, a group of order $48$. If $R \to R'$ is flat, then Zariski locally on $R$ we may apply (\ref{prop:A7.1.3}, (1)), by inverting an odd prime $\ell$ and viewing $\mathcal{P}$ as the quotient of $(\mathcal{P}, [\Gamma(\ell)])$ by the group $\GL(2, \F_\ell)$.
\end{proof}

\begin{krem}\label{rem:8.1.7}
Formation of the coarse moduli scheme does not always commute with base change.
\end{krem}

\begin{pitemnum}{8.1.7.1}\label{item:8.1.7.1}
The simplest example is provided by the moduli problem $[\omega]$ on \EllZ\ defined by
\[
\begin{aligned}
[\omega](E/S) =\ &\text{the set of nowhere vanishing invariant one-forms} \\
&\text{on $E/S$ .}
\end{aligned}
\]
This $[\omega]$ is a $\mathbb{G}_m$-torsor on \EllZ. For any ring $R$, the coarse moduli scheme $\Mcoarse([\omega] \otimes R)$ is the graded ring $\MF(R)$ of all level one modular forms (with no condition at infinity) over $R$. It is known that (cf.\ \cite{De2})
\[
\begin{aligned}
\MF(\Z) &= \Z[c_4, c_6, \Delta, \Delta^{-1}]/(c_4^3 - c_6^2 = 1728\, \Delta) \\
\MF(\F_3) &= \F_3[a_2, \Delta, \Delta^{-1}] \ ; \quad a_2 = \text{Hasse invar.} \\
\MF(\F_2) &= \F_2[a_1, \Delta, \Delta^{-1}] \ ; \quad a_1 = \text{Hasse invar.} ,
\end{aligned}
\]
with interrelations
\[
\left\{
\begin{array}{ll}
c_4 \to a_2^2 , \quad c_6 \to a_2^3 & \text{mod } 3 \\[2pt]
c_4 \to a_1^4 , \quad c_6 \to a_1^6 & \text{mod } 2 .
\end{array}
\right.
\]
Thus the maps
\[
\left\{
\begin{array}{l}
\MF(\Z) \otimes \F_3 \to \MF(\F_3) \\[2pt]
\MF(\Z) \otimes \F_2 \to \MF(\F_2)
\end{array}
\right.
\]
both fail to be isomorphisms.
\end{pitemnum}

\begin{pitemnum}{8.1.7.2}\label{item:8.1.7.2}
For an example which is finite flat over \EllZ, consider the sub-problem $[\Delta = 1]$ defined by
\[
\begin{aligned}
[\Delta = 1](E/S) =\ &\text{the set of nowhere vanishing translation-} \\
&\text{invariant $\omega$'s such that $\Delta(E, \omega) = 1$ .}
\end{aligned}
\]
This $[\Delta = 1]$ is a $\mu_{12}$-torsor over \EllZ. Over any ring $R$, it is easy to see that
\[
\Mcoarse([\Delta = 1] \otimes R) = \Spec\bigl(\MF(R)/(\Delta - 1)\bigr) ,
\]
and so
\[
\begin{aligned}
\Mcoarse([\Delta = 1]) &= \Spec\bigl(\Z[c_4, c_6]/(c_4^3 - c_6^2 = 1728)\bigr) \\
\Mcoarse([\Delta = 1] \otimes \F_3) &= \Spec\bigl(\F_3[a_2]\bigr) \\
\Mcoarse([\Delta = 1] \otimes \F_2) &= \Spec\bigl(\F_2[a_1]\bigr) .
\end{aligned}
\]
Once again, base change fails at the primes two and three.
\end{pitemnum}

% ------------------------------------------------------------------------------
% Section 8.2
% ------------------------------------------------------------------------------
\section{The $j$-line as a coarse moduli scheme}
\label{sec:8.2}

\kmitem\label{item:8.2.1}
It is well known (cf.\ \cite[\S2]{Ig3}) that over any ring $R$, the coarse moduli scheme attached to the moduli problem $[\Gamma(1)]$ is the $j$-line over $R$
\[
\Mcoarse([\Gamma(1)]) = \Spec(R[j]) ,
\]
with the $j$-invariant normalized \`a la Tate (cf.\ \cite{De2}); $j = 0$ has complex multiplication by $\Z[\zeta_3]$, $j = 1728$ by $\Z[i]$.

For any relatively representable $\mathcal{P}$, affine over \EllR, the unique map of moduli problems on \EllR
\[
\mathcal{P} \to [\Gamma(1)]
\]
defines a morphism of $R$-schemes
\[
\Mcoarse(\mathcal{P}) \to \Spec(R[j]) .
\]

\begin{kprop}\label{prop:8.2.2}
If $R$ is noetherian, and if $\mathcal{P}$ is finite over \EllR, then
\[
\Mcoarse(\mathcal{P}) \to \Spec(R[j])
\]
is finite.
\end{kprop}

\begin{proof}
The question is local on $R$, so we may assume some odd prime $\ell$ is invertible in $R$. We have a commutative diagram of $R$-schemes
\[
\begin{tikzcd}[column sep=large, row sep=large]
& \Mfrak(\mathcal{P}, [\Gamma(\ell)]) \arrow[ld, "{\text{finite by hypo.}}"'] \arrow[rd, "{\text{divide by } \GL(2,\F_\ell)}"] & \\
\Mfrak([\Gamma(\ell)]) \arrow[rd, "{\text{divide by } \GL(2,\F_\ell)}"'] & & \Mfrak(\mathcal{P}) \arrow[ld] \\
& \Mcoarse([\Gamma(1)]) &
\end{tikzcd}
\]
Because $\Mfrak([\Gamma(\ell)])$ is an affine $R$-scheme of finite type, it is finite over its quotient $\Mcoarse([\Gamma(1)])$ by the finite group $\GL(2, \F_\ell)$. Therefore $\Mfrak(\mathcal{P}, [\Gamma(\ell)])$ is finite over $\Mcoarse([\Gamma(1)])$. Because $R$ is noetherian, $\Mcoarse([\Gamma(1)])$ is the Spec of a noetherian ring, and therefore the coordinate ring of $\Mcoarse(\mathcal{P})$, being contained in a finite $R[j]$-module, is itself finite over $R[j]$.
\end{proof}

\begin{kprop}\label{prop:8.2.3}
Let $\mathcal{P}$ be relatively representable and finite over \EllR. Let $B$ be a complete local noetherian $R$-algebra which is flat over $R$, and whose residue field $k$ is algebraically closed. Let $E/k$ be an elliptic curve, and denote by $E/B[[T]]$ its universal formal deformation (to artin local $B$-algebras). Let $j_{00} \in k$ be the $j$-invariant $j(E/k)$, and $j_0 \in B$ some lifting of it to $B$. Then we have a cartesian diagram
\[
\begin{tikzcd}[column sep=huge, row sep=large]
\Mcoarse(\mathcal{P}) & (\mathcal{P}_{E/B[[T]]})/\Aut(E/k) \arrow[l] \arrow[d] \\
\Spec(R[j]) \arrow[u] & \Spec(B[[j - j_0]]) \arrow[l]
\end{tikzcd}
\]
\end{kprop}

\begin{proof}
Because $B$ is $R$-flat, $\Mcoarse(\mathcal{P} \tens{R} B) \xrightarrow{\ \sim\ } \Mcoarse(\mathcal{P}) \tens{R} B$, and we are reduced to the case $R = B$. The action of $\Aut(E/k)$ on the finite $B[[T]]$-scheme $\mathcal{P}_{E/B[[T]]}$ is by viewing it as the formal $B$-scheme which pro-represents the functor on artin local $B$-algebras with residue field $k$
\[
\begin{aligned}
A \mapsto\ &\text{triples } (\widetilde{E}/A, a, i) \text{ where } \widetilde{E}/A \text{ is an elliptic curve,} \\
&\text{$a \in \mathcal{P}(\widetilde{E}/A)$, and $i$ is a $k$-isomorphism } \widetilde{E} \tens{R} k \xrightarrow{\ \sim\ } E .
\end{aligned}
\]
The group $\Aut(E/k)$ acts on this functor by moving the isomorphism $i$.

Choose an odd prime $\ell \ne \chark(k)$. Then $\ell$ is invertible in $B$, so that $\mathcal{S} = [\Gamma(\ell)] \tens{\Z} B$ is a finite etale $\GL(2, \F_\ell)$-torsor over $(\mathrm{Ell}/B)$. By definition, we have
\[
\Mcoarse(\mathcal{P}) \xrightarrow{\ \sim\ } \Mfrak(\mathcal{P}, \mathcal{S})/\GL(2, \F_\ell) ,
\]
as $B[j]$-schemes. By flatness of $B[[j - j_0]]$ over $B[j]$, we infer
\[
\Mcoarse(\mathcal{P}) \tens{B[j]} B[[j - j_0]] \xrightarrow{\ \sim\ } \bigl(\Mfrak(\mathcal{P}, \mathcal{S}) \tens{B[j]} B[[j - j_0]]\bigr)/\GL(2, \F_\ell) .
\]
If we knew the proposition for $(\mathcal{P}, \mathcal{S})$, we could rewrite this
\[
\begin{aligned}
&= \bigl((\mathcal{P}, \mathcal{S})_{E/B[[T]]}\bigr)/\Aut(E/k) \times \GL(2, \F_\ell) \\
&= \bigl((\mathcal{P}, \mathcal{S})_{E/B[[T]]}/\GL(2, \F_\ell)\bigr)/\Aut(E/k) \\
&= \bigl(\mathcal{P}_{E/B[[T]]}\bigr)/\Aut(E/k) ,
\end{aligned}
\]
the last equality because $(\mathcal{P}, \mathcal{S})$ is a finite etale $\GL(2, \F_\ell)$-torsor over $\mathcal{P}$, so we can pass to the quotient ``object by object.''

Thus we are reduced to the case when $\mathcal{P}$ is representable and finite over $(\mathrm{Ell}/B)$. Then $\Mfrak(\mathcal{P})$ is finite over $B[j]$, so we have a canonical decomposition, $B$ being complete,
\[
\Mfrak(\mathcal{P}) \tens{B[j]} B[[j - j_0]] = \coprod_{j(x)=j_0} \Spec\bigl(\widehat{\cO}_{\Mfrak(\mathcal{P}), x}\bigr)
\]
over the $k$-valued (= closed) points of $\Mfrak(\mathcal{P})$ with $j$-invariant $j_0$. Each such point $x$ is of the form $(E/k$, some $a_x \in \mathcal{P}(E/k))$, and the $a_x$ run over a set of representatives of $\mathcal{P}(E/k)$ modulo $\Aut(E/k)$. By definition of $\Mfrak(\mathcal{P})$, the corresponding complete local ring $\widehat{\cO}_{\Mfrak(\mathcal{P}), x}$ pro-represents the functor on artin local $B$-algebras with residue field $k$
\[
A \mapsto
\left\{
\begin{array}{l}
\text{pairs } (\widetilde{E}/A, a \in \mathcal{P}(\widetilde{E}/A)) \text{ such that} \\[2pt]
(\widetilde{E}/A, a) \tens{A} k \text{ admits an isomorphism} \\[2pt]
\text{with } (E/k, a_x) .
\end{array}
\right.
\]
Because $\mathcal{P}$ is representable, it is rigid, so $\Aut(E/k)$ acts freely on $\mathcal{P}(\widetilde{E}/k)$; therefore if $(\widetilde{E}/A, a) \tens{A} k$ admits an isomorphism with $(E/k, a_x)$, this isomorphism is unique.

Therefore the finite $B[[T]]$-scheme
\[
\Mfrak(\mathcal{P}) \tens{B[j]} B[[j - j_0]]
\]
pro-represents the functor
\[
A \mapsto
\left\{
\begin{array}{l}
\text{an element } a_x \in \mathcal{P}(E/k) \text{ among our chosen} \\[2pt]
\text{set of representatives, together with a triple} \\[2pt]
(\widetilde{E}/A, a \in \mathcal{P}(\widetilde{E}/A), i) \text{ with } i \text{ an isomorphism} \\[2pt]
(\widetilde{E}/A, a) \tens{A} k \xrightarrow{\ \sim\ } (E/k, a_x) .
\end{array}
\right.
\]
On the other hand, $\mathcal{P}_{E/B[[T]]}$ pro-represents the functor
\[
A \mapsto
\left\{
\begin{array}{l}
\text{an element } a_0 \in \mathcal{P}(E/k) \text{ and a triple} \\[2pt]
(\widetilde{E}/A, a \in \mathcal{P}(\widetilde{E}/A), i) \text{ where } i \text{ is an} \\[2pt]
\text{isomorphism} \\[2pt]
(\widetilde{E}/A, a) \tens{A} k \xrightarrow{\ \sim\ } (E/k, a_0) .
\end{array}
\right.
\]
By the rigidity of $\mathcal{P}$, the isomorphism $i$ is unique. As $\Aut(E/k)$ operates freely on the indexing set $\mathcal{P}(E/k)$ of $a_0$'s we may identify the quotient of $\mathcal{P}_{E/B[[T]]}$ by $\Aut(E/k)$ to the subfunctor of $\mathcal{P}_{E/B[[T]]}$ defined by requiring $a_0$ to lie in a fixed set of representatives, say the $a_x$'s, of $\mathcal{P}(E/k)$ modulo $\Aut(E/k)$.
\end{proof}

% ------------------------------------------------------------------------------
% Section 8.3
% ------------------------------------------------------------------------------
\section{Localization of moduli problems over the $j$-line}
\label{sec:8.3}

\kmitem\label{item:8.3.1}
Let $U$ be an open set in the affine $j$-line $\Spec(R[j])$. We denote by $(\mathrm{Ell}/U/R)$ the full subcategory of \EllR\ consisting of those elliptic curves $E/S/R$ whose $j$-invariant ``lies in $U$'', in the sense that the morphism ``$j$'' factors through $U$ :
\begin{equation}\label{eq:8.3.1.1}
\begin{tikzcd}[column sep=huge, row sep=large]
S \arrow[d, "j"] \arrow[rd, dashed] & \\
\Spec(R[j]) \supset & U .
\end{tikzcd}
\end{equation}
We denote by $[U]$ the moduli problem on \EllR\ which is the ``characteristic function of $U$''; for any $E/S/R$, we define
\begin{equation}\label{eq:8.3.1.2}
[U](E/S) = \left\{
\begin{array}{l}
\text{the set with one element if} \\
\text{$E/S/R$ lies in $(\mathrm{Ell}/U/R)$,} \\
\text{the empty set if not.}
\end{array}
\right.
\end{equation}
This moduli problem is relatively representable and etale (indeed an open immersion) over \EllR; for any $E/S/R$, we have
\begin{equation}\label{eq:8.3.1.3}
[U]_{E/S} = S \fib{\Spec(R[j])} U .
\end{equation}

\kmitem\label{item:8.3.2}
For any moduli problem $\mathcal{P}$ on \EllR, we define $\mathcal{P}|U$ to be the open sub-problem defined by
\begin{equation}\label{eq:8.3.2.1}
\begin{aligned}
(\mathcal{P}|U)(E/S) &= \mathcal{P}(E/S) \times [U](E/S) \\
&= \left\{
\begin{array}{l}
\mathcal{P}(E/S) \text{ if } E/S \text{ lies in } (\mathrm{Ell}/U/R) \\
\text{the empty set if not.}
\end{array}
\right.
\end{aligned}
\end{equation}
If $\mathcal{P}$ is relatively representable over \EllR, then so is $\mathcal{P}|U$, and we have
\begin{equation}\label{eq:8.3.2.2}
(\mathcal{P}|U)_{E/S} = \mathcal{P}_{E/S} \fib{\Spec(R[j])} U \quad \underline{\text{notation}} \quad \mathcal{P}_{E/S}|U .
\end{equation}
In particular, if $\mathcal{P}$ is representable, then $\mathcal{P}|U$ is representable and we have
\begin{equation}\label{eq:8.3.2.3}
\Mfrak(\mathcal{P}|U) = \Mfrak(\mathcal{P}) \fib{\Spec(R[j])} U = \Mfrak(\mathcal{P})|U .
\end{equation}
If $U$ is affine over $R$, then for any relatively representable $\mathcal{P}$ which is affine over \EllR, $\mathcal{P}|U$ is affine over \EllR, and the coarse moduli schemes of $\mathcal{P}$ and of $\mathcal{P}|U$ are related by
\begin{equation}\label{eq:8.3.2.4}
\Mcoarse(\mathcal{P}|U) = \Mcoarse(\mathcal{P}) \fib{\Spec(R[j])} U = \Mcoarse(\mathcal{P})|U .
\end{equation}
In particular, taking $\mathcal{P} = [\Gamma(1)]$, we find
\begin{equation}\label{eq:8.3.2.5}
\Mcoarse([U]) = U .
\end{equation}

% ------------------------------------------------------------------------------
% Section 8.4
% ------------------------------------------------------------------------------
\section{The $j$-invariant as a fine modulus, coarse moduli schemes as fine moduli schemes (!)}
\label{sec:8.4}

\kmitem\label{item:8.4.1}
Recall from \cite{De2} that the ring of modular forms over $\Z$ is generated by the forms $c_4$, $c_6$, $\Delta$, and $\Delta^{-1}$, with the single relation
\[
(c_4)^3 - (c_6)^2 = 1728\Delta
\]
and that $j$ is given by the formulas
\begin{gather*}
j = (c_4)^3/\Delta \\
j - 1728 = (c_6)^2/\Delta .
\end{gather*}

\kmitem\label{item:8.4.2}
Recall also that for an elliptic curve $E/S$, the following conditions are equivalent
\begin{enumerate}[(1)]
\item For any connected $S$-scheme $T$, the elliptic curve $E_T/T$ has $\Aut(E_T/T) = \pm 1$.
\item $j(E/S)(j(E/S) - 1728)$ is invertible in $\Gamma(S, \cO_S)$.
\item For any $S$-scheme $T$, and any nowhere vanishing invariant differential $\omega$ on $E_T/T$, the elements $c_4(E_T/T, \omega)$, $c_6(E_T/T, \omega)$ are both invertible in $\Gamma(T, \cO_T)$.
\item $E/S$ lies in $(\mathrm{Ell}/U/\Z)$ for $U$ the open set $\Spec(\Z[j][1/(j(j-1728))])$ of $\Spec(\Z[j])$.
\end{enumerate}

One also has the following more precise result, which we quote from \cite{De2}.

\begin{kprop}\label{prop:8.4.3}
Let $S$ be an arbitrary scheme, $E_1/S$ and $E_2/S$ two elliptic curves over $S$. Suppose that $j(E_1/S) = j(E_2/S)$, and that $j(E_1/S) \cdot (j(E_1/S) - 1728)$ is invertible in $\Gamma(S, \cO_S)$. Then the functor on $(\Sch/S)$
\[
\begin{aligned}
T \mapsto\ &\text{isomorphisms of elliptic curves over } T \\
&\text{from } E_1 \fib{S} T \text{ to } E_2 \fib{S} T ,
\end{aligned}
\]
is represented by an $S$-scheme $\Isom_S(E_1, E_2)$ which is a finite etale galois covering of $S$ with group $\pm 1$.
\end{kprop}

\begin{kcor}\label{cor:8.4.4}
Let $R$ be an arbitrary ring, $\mathcal{P}$ and $\mathcal{S}$ two representable moduli problems on \EllR.
Let $(\mathcal{P}, \mathcal{S})$ denote the simultaneous moduli problem
\[
E/S \mapsto \mathcal{P}(E/S) \times \mathcal{S}(E/S) ,
\]
which sits in a commutative diagram of moduli problems on \EllR
\[
\begin{tikzcd}[column sep=huge, row sep=large]
& (\mathcal{P}, \mathcal{S}) \arrow[ld] \arrow[rd] & \\
\mathcal{P} \arrow[rd] & & \mathcal{S} \arrow[ld] \\
& {[\Gamma(1)]} &
\end{tikzcd}
\]
The induced morphisms of coarse moduli schemes
\[
\begin{tikzcd}[column sep=huge, row sep=large]
& \Mfrak(\mathcal{P}, \mathcal{S}) \arrow[ld] \arrow[rd] & \\
\Mfrak(\mathcal{P}) \arrow[rd, "j"] & & \Mfrak(\mathcal{S}) \arrow[ld, "j"'] \\
& \Aff^1_R &
\end{tikzcd}
\]
define a morphism of $\Aff^1_R$-schemes
\[
\Mfrak(\mathcal{P}, \mathcal{S}) \longrightarrow \Mfrak(\mathcal{P}) \fib{\Aff^1_R} \Mfrak(\mathcal{S}) .
\]
This morphism, over the open set $U$ in $\Aff^1_R$ where $j(j-1728)$ in invertible, makes $\Mfrak(\mathcal{P}, \mathcal{S})|U$ into a finite etale galois covering of $\Mfrak(\mathcal{P})|U \fib{U} \Mfrak(\mathcal{S})|U$, with group $\pm 1$, and so it induces an isomorphism of $U$-schemes
\[
(\Mfrak(\mathcal{P}, \mathcal{S})|U)/{\pm 1} \cong \Mfrak(\mathcal{P})|U \fib{U} \Mfrak(\mathcal{S})|U .
\]
\end{kcor}

\begin{proof}
Consider the functor $Z$ on $(\Sch/R)$ represented by $\Mfrak(\mathcal{P})|U \fib{U} \Mfrak(\mathcal{S})|U$. For any $R$-scheme $S$, we have
\[
\begin{aligned}
Z(S) =\ &\text{$S$-isomorphism classes of data } (E_1/S, \alpha, E_2/S, \beta) \\
&\text{with } E_1/S, E_2/S \text{ elliptic curves over } S \text{ with} \\
&\text{$j(E_1/S) = j(E_2/S)$, $j(E_1/S)(j(E_1/S) - 1728)$} \\
&\text{invertible in } \Gamma(S, \cO_S), \alpha \in \mathcal{P}(E_1/S), \beta \in \mathcal{S}(E_2/S).
\end{aligned}
\]
Now consider the functor $W$ on $(\Sch/R)$ represented by $\Mfrak(\mathcal{P}, \mathcal{S})|U$. For any $R$-scheme $S$, we have
\[
\begin{aligned}
W(S) =\ &\text{$S$-isomorphism classes of data } (E/S, \alpha, \beta) \\
&\text{with } E/S \text{ an elliptic curve over } S \text{ with} \\
&\text{$j(E/S)(j(E/S) - 1728)$ invertible in } S, \\
&\alpha \in \mathcal{P}(E/S), \beta \in \mathcal{S}(E/S).
\end{aligned}
\]
The morphism of $U$-schemes
\[
\Mfrak(\mathcal{P}, \mathcal{S})|U \to \Mfrak(\mathcal{P})|U \fib{U} \Mfrak(\mathcal{S})|U
\]
corresponds to the map on functors
\[
W \to Z
\]
defined on $S$-valued points by
\[
(E/S, \alpha, \beta) \mapsto (E/S, \alpha, E/S, \beta) .
\]
Therefore if we are given any $R$-morphism $S \to \Mfrak(\mathcal{P})|U \fib{U} \Mfrak(\mathcal{S})|U$, corresponding to data $(E_1/S, \alpha, E_2/S, \beta)$, then the fiber product
\[
\begin{tikzcd}[column sep=large, row sep=large]
\square \arrow[r] \arrow[d] & S \arrow[d] \\
\Mfrak(\mathcal{P}, \mathcal{S})|U \arrow[r] & \Mfrak(\mathcal{P})|U \fib{U} \Mfrak(\mathcal{S})|U ,
\end{tikzcd}
\]
is none other than the scheme
\[
\Isom_S(E_1, E_2) ,
\]
which is finite etale galois over $S$ with group $\pm 1$.
\end{proof}

\begin{kcor}\label{cor:8.4.5}
Let $R$ be a regular noetherian ring, $\mathcal{P}$ a representable moduli problem on \EllR. Let $U$ be an affine open set in $\Spec(R[j])$. Suppose that $j(j-1728)$ is invertible on $U$, and suppose that $\mathcal{P}|U$ is finite etale over $[U]$. Then the morphism ``$j$''
\[
\Mfrak(\mathcal{P}) \to \Spec(R[j])
\]
is finite etale over $U$, i.e., the map
\[
\Mfrak(\mathcal{P}|U) \to U
\]
is finite etale.
\end{kcor}

\begin{proof}
Because $R$ is regular, and $\mathcal{P}|U$ is finite etale over $[U]$, the moduli scheme $\Mfrak(\mathcal{P}|U)$ is regular, as is $U$ itself. The morphism $\Mfrak(\mathcal{P}|U) \to U$ is certainly finite, so it is automatically flat, because source and target are regular schemes of the same dimension. Because $\Mfrak(\mathcal{P}|U) \to U$ is finite and flat, we may base-change by it to test whether some other map $X \to U$ is finite etale. Thus it suffices to see that
\[
\Mfrak(\mathcal{P}|U) \fib{U} \Mfrak(\mathcal{P}|U) \xrightarrow{\ \mathrm{pr}_1\ } \Mfrak(\mathcal{P}|U)
\]
is finite etale. This map sits under
\[
\Mfrak(\mathcal{P}|U, \mathcal{P}|U) \longrightarrow \Mfrak(\mathcal{P}|U) \fib{U} \Mfrak(\mathcal{P}|U) \xrightarrow{\ \mathrm{pr}_1\ } \Mfrak(\mathcal{P}|U) .
\]
The first map is a finite etale $\pm 1$ torsor, [$\pm 1$ acting on pairs of level $\mathcal{P}$ structures by $(a_1, a_2) \mapsto (a_1, \pm a_2)$], and the composite
\[
\Mfrak(\mathcal{P}|U, \mathcal{P}|U) \to \Mfrak(\mathcal{P}|U)
\]
is finite etale by hypothesis. Therefore our map is also finite etale.
\end{proof}

\begin{kprop}\label{prop:8.4.6}
Let $R$ be any ring, $U$ any affine open set of $\Spec(R[j])$ on which $j(j-1728)$ is invertible, $\mathcal{P}$ a relatively representable moduli problem which is affine over \EllR. Suppose that for any $E/S/R$ which lies in $(\mathrm{Ell}/U/R)$, the automorphism $[-1]_E$ operates on the set $\mathcal{P}(E/S)$ without fixed points. Then $\mathcal{P}|U$ is representable, and
\[
\Mfrak(\mathcal{P}|U) = \Mcoarse(\mathcal{P}|U) = \Mcoarse(\mathcal{P})|U .
\]
\end{kprop}

\begin{proof}
The hypothesis assures that $\mathcal{P}|U$ is rigid.
\end{proof}

\begin{kcor}\label{cor:8.4.7}
Let $N \ge 1$ be any integer. Let $U_N$ be an affine open set in $\Spec(\Z[j])$ such that
\[
\begin{cases}
\text{$j(j-1728)$ is invertible on $U_N$} \\
\text{for all $p|N$, $U_N(\overline{\F}_p)$ contains no $j$'s of} \\
\text{supersingular curves in characteristic $p$ .}
\end{cases}
\]
Then
\begin{enumerate}[(1)]
\item If $N \ge 3$, then $[\Gamma(N)]|U_N$ is representable on $(\mathrm{Ell}/\Z)$.
\item If $N = 1,2,4$, then $[\mathrm{bal.}\ \Gamma_1(N)]|U_N$ is representable on $(\mathrm{Ell}/\Z)$.
\end{enumerate}
\end{kcor}

\begin{proof}
Of course if $N$ is divisible by two relatively prime integers both $\ge 3$, (resp.\ both $\ge 4$), then $[\Gamma(N)]$ (resp.\ $[\mathrm{bal.}\ \Gamma_1(N)]$) itself is representable.

To prove (1), it suffices to show that if $E/S$ is an elliptic curve with a $\Gamma(N)$-structure $(P, Q)$ with $2P = 2Q = 0$, then there is a geometric point $s$ of $S$ of characteristic $p|N$ at which $E_s$ is supersingular. If $N$ is odd, $2P = 2Q = 0$ implies $P = Q = 0$, whence for any prime $p|N$, $((N/p)P, (N/p)Q) = (0,0)$ is a Drinfeld $p$-basis for $E/S$ (\ref{thm:5.5.7}). Therefore $S$ is an $\F_p$-scheme, and $E_s$ is supersingular at every geometric point of $S$ (\ref{lem:5.3.2.1}, \ref{lem:5.3.6}). In particular, $N$ is a power of $p$. If $N$ is even, then $(2P, 2Q) = (0,0)$ is a Drinfeld $N/2$-structure in $E/S$, and we repeat the same argument with $N/2$, which by hypothesis is still $\ge 2$.

To prove (2), it suffices to show that if $(E \to E', P, P')$ is a balanced $\Gamma_1(N)$ structure over some $S$, with $2P = 2P' = 0$, then there is a geometric point $s$ of $S$ of characteristic $p|N$ at which $E_s$ is supersingular. If $N$ is odd, $2P = 2P' = 0$ implies $P = P' = 0$. Therefore (\ref{cor:5.5.5}) $(0,0)$ is a $\Gamma(N)$-structure on $E/S$, and we conclude as before.

If $N$ is even, we argue as follows. First suppose $N = 2^k M$ with $M$ odd $\ge 3$. Consider the standard factorization of our given cyclic $E \to E'$ into a cyclic $M$-isogeny followed by a cyclic $2^k$-isogeny:
\[
P;\; E \xleftrightarrow[\ \pi^t_{\mathrm{odd}}\ ]{\ \pi_{\mathrm{odd}}\ } E'' \xleftrightarrow[\ \pi^t_{\mathrm{even}}\ ]{\ \pi_{\mathrm{even}}\ } E';\; P' .
\]
By the Backing-up Theorem (\ref{thm:6.7.11}), $2^k P$ generates $\Ker(\pi_{\mathrm{odd}})$ and $\pi^t_{\mathrm{even}}(P')$ generates $\Ker(\pi^t_{\mathrm{odd}})$. As $2P = 2P' = 0$, the previous case ``$N$ odd'' applies to
\[
2^k P;\; E \xleftrightarrow[\ \pi^t_{\mathrm{odd}}\ ]{\ \pi_{\mathrm{odd}}\ } E'';\; \pi^t_{\mathrm{even}}(P') .
\]
If $N = 2^n$, we consider the standard factorization into a sequence of 2-isogenies
\[
P = P_0;\; E = E_0 \xleftrightarrow[\ \pi^t_{0,1}\ ]{\ \pi_{0,1}\ } E_1 \cdots E_{n-1} \xleftrightarrow[\ \lambda_{n,n-1}\ ]{\ \pi_{n-1,n}\ } E_n = E';\; P' = Q_n .
\]
Then $2\pi_{0,1}(P_0)$ generates $\Ker(\pi_{1,n-1})$, and $2\lambda_{n,n-1}(Q_n)$ generates $\Ker(\lambda_{n-1,1})$. As $n$ is $\ge 3$, these are dual $2^{n-2}$-isogenies, both generated by zero, so by (\ref{lem:5.3.2.3}) $S$ is an $\F_2$-scheme and $E_1/S$ is supersingular at every geometric point of $S$. Since $E_0$ is isogenous to $E_1$, it too is supersingular at every geometric point of $S$.
\end{proof}

\begin{krem}\label{rem:8.4.8}
The moduli problems $[\Gamma(1)]$, $[\Gamma(2)]$, $[\mathrm{bal.}\ \Gamma_1(2)]$ are not representable, over any non-void open set of the $j$-line, because they are not \emph{rigid}; the automorphism $[-1]$ visibly acts trivially. As for $[\mathrm{bal.}\ \Gamma_1(4)]$, it is of course representable over $\Z[1/2]$, but it is not representable over \emph{any} non-void open set of the $j$-line which contains any point of characteristic two. For such an open meets the $j$-line over $\F_2$ in a non-void open set, hence contains ordinary points in characteristic two. If $E/k$ is an ordinary curve over an algebraically closed field $k$ of characteristic $2$, then there exists an isomorphism
\[
E[4] \cong \mu_4 \times \Z/4\Z .
\]
The points
\begin{gather*}
P = (1,2) \\
Q = (1,1)
\end{gather*}
form a $\Gamma(4)$-structure (because $P - 2Q$ generates the $\mu_4$, and $Q$ generates the $\Z/4\Z$). Visibly $2Q = P$, so the associated balanced $\Gamma_1(4)$-structure $(P, Q \bmod P)$ is fixed by the automorphism $[-1]$.
\end{krem}

\begin{kthm}\label{thm:8.4.9}
Let $R$ be a ring in which $2$ is invertible, and let $a \in R^\times$ be a unit of $R$. Let $U \subset \Spec(R[j])$ be the open set $U = \Spec(R[j][1/(j(j-1728))])$ where $j(j-1728)$ is invertible. Denote by $[U, a]$ the moduli problem on \EllR\ defined by
\[
\begin{aligned}
[U, a](E/S) =\ &\text{the set of nowhere vanishing invariant differentials} \\
&\text{$\omega$ on $E/S$ such that $c_4(E, \omega)$ and $c_6(E, \omega)$ are} \\
&\text{both invertible on $S$, and such that} \\[3pt]
&c_6(E, \omega)/c_4(E, \omega) = a .
\end{aligned}
\]
Then we have
\begin{enumerate}[(1)]
\item The natural morphism ``forget $\omega$'' of moduli problems on \EllR
\[
[U, a] \mapsto [U]
\]
is finite etale and galois with galois group $\pm 1$.
\item The moduli problem $[U, a]$ is representable, and is affine and etale over \EllR.
\item The map of coarse moduli schemes induced by (1),
\[
\Mfrak([U, a]) \to \Mcoarse([U]) = U
\]
is an isomorphism of $R$-schemes.
\end{enumerate}
\end{kthm}

\begin{proof}
It is clear from the formulas expressing $j$ via $c_4$, $c_6$ and $\Delta$, that the natural map $[U, a] \to [\Gamma(1)]$ factors through $[U]$. To show that (1) holds, consider an $E/S/R$ such that its $j(j-1728)$ is invertible on $S$. Zariski localizing on $S$, we may suppose there exists a nowhere-vanishing invariant differential $\omega_0$ on $E/S$. Then both $c_4(E, \omega_0)$ and $c_6(E, \omega_0)$ are invertible on $S$. Then a nowhere-vanishing $\omega$ on $E/S$ may be written in the form $\lambda^{-1}\omega_0$ for some unit $\lambda \in \Gamma(S, \cO_S)^\times$, and such an $\omega$ lies in $[U, a](E/S)$ if and only if
\[
\lambda^2 = a \cdot \frac{c_4(E, \omega_0)}{c_6(E, \omega_0)} .
\]
Because $2$ is invertible, this equation defines $[U, a]_{E/S}$ as a finite etale galois covering of $S$ with group $\pm 1$. Thus (1) holds, and, $[U]$ being affine and etale over \EllR, we see that its finite etale covering $[U, a]$ is itself relatively representable, affine and etale over \EllR. Because $[U, a]$ is a \emph{rigid} moduli problem, it is consequently representable, necessarily by a smooth affine curve $\Mfrak([U, a])$ over $R$.

Because $[U]$ is the quotient of $[U, a]$ by the group $\pm 1$, we have
\[
\Mfrak([U, a])/{\pm 1} \cong \Mcoarse([U]) = U .
\]
But the group $\pm 1$ acts trivially on the scheme $\Mfrak([U, a])$, because as an $R$-scheme $\Mfrak([U, a])$ represents the functor on $R$-schemes
\[
\begin{aligned}
S \mapsto\ &\text{$S$-isomorphism classes of pairs } (E/S, \omega) \\
&\text{where $\omega$ is a $[U, a]$-structure on } E/S ,
\end{aligned}
\]
and the data $(E/S, \omega)$ is $S$-isomorphic (by $[-1]_E$) to the data $(E/S, -\omega)$.
\end{proof}

\begin{kcor}\label{cor:8.4.10}
Let $R$ be a ring in which $2$ is invertible, $a \in R^\times$ a unit, $\mathcal{P}$ a relatively representable moduli problem which is affine over \EllR. Let $U$ denote the open set of $\Spec(R[j])$ where $j(j-1728)$ is invertible. Then
\begin{enumerate}[(1)]
\item The simultaneous moduli problem $(\mathcal{P}, [U, a])$ is representable, and it is finite etale galois over $\mathcal{P}|U$ with galois group $\pm 1$.
\item We have a natural isomorphism of $R$-schemes
\[
\Mfrak(\mathcal{P}, [U, a])/{\pm 1} \cong \Mcoarse(\mathcal{P}|U) = \Mcoarse(\mathcal{P})|U .
\]
\item If, for every elliptic curve $E/S/R$, the automorphism $[-1]_E$ acts trivially on the set $\mathcal{P}(E/S)$, then the isomorphism in (2) above is an isomorphism of $R$-schemes
\[
\Mfrak(\mathcal{P}, [U, a]) \cong \Mcoarse(\mathcal{P}|U) = \Mcoarse(\mathcal{P})|U .
\]
\end{enumerate}
\end{kcor}

\begin{proof}
Assertion (1) holds because $[U, a]$ is representable, and finite etale galois over $[U]$ with group $\pm 1$. By (\ref{lem:8.1.5}), this implies (2). Under the hypothesis of (3), the group $\pm 1$ acts trivially on $\Mfrak(\mathcal{P}, [U, a])$.
\end{proof}

\begin{kthm}\label{thm:8.4.11}
Let $R$ be a ring in which $2$ is nilpotent, $U$ the open set of $\Spec(R[j])$ where $j(j-1728)$ is invertible. Denote by $[\mu_4]$ the moduli problem on \EllR\ defined by
\[
[\mu_4](E/S) = \text{the set of $S$-group injections } \mu_4 \xhookrightarrow{\ a\ } E .
\]
Then
\begin{enumerate}[(1)]
\item The natural morphism ``forget $a$'' of moduli problems on \EllR
\[
[\mu_4] \to [\Gamma(1)]
\]
factors through $[U]$, and the induced map
\[
[\mu_4] \to [U]
\]
is finite etale and galois with galois group $\pm 1$.
\item The moduli problem $[\mu_4]$ is representable, and it is affine and etale over \EllR.
\item The map of coarse moduli schemes induced by (1)
\[
\Mfrak([\mu_4]) \to \Mcoarse([U]) = U
\]
is an isomorphism of $R$-schemes.
\end{enumerate}
\end{kthm}

\begin{proof}
Because $2$ is nilpotent, an $E/S/R$ has $j(j-1728)$ invertible if and only if its $j$ is invertible if and only if $E/S$ is fiber-by-fiber ordinary.

Now when $2$ is nilpotent, $E/S$ is fiber-by-fiber ordinary if and only if locally f.p.p.f. on $S$, the kernel of $4$ in the formal group $\widehat{E}[4]$ is isomorphic to $\mu_4$, if and only if locally f.p.p.f. on $S$ there exists an $S$-group inclusion $\mu_4 \xhookrightarrow{\ a\ } E$. Therefore $[\mu_4]$ does indeed lie over $[U]$, and over $[U]$ it is the finite etale galois covering with galois group $\pm 1 = (\Z/4\Z)^\times$ defined by
\[
[\mu_4]_{E/S} = \Isom_S(\mu_4, \widehat{E}[4])
\]
for any fiber-by-fiber ordinary $E/S/R$. Assertion (2) is in fact true over $\Z$, as has already been noted, because $[\mu_4]$ is a rigid problem. Assertion (3) holds because $[U] = [\mu_4]/{\pm 1}$, while $\pm 1$ acts trivially on the scheme $\Mfrak([\mu_4])$.
\end{proof}

\begin{kcor}\label{cor:8.4.12}
Let $R$ be a ring in which $2$ is nilpotent, $U$ the open set of $\Spec(R[j])$ where $j(j-1728)$ is invertible, $\mathcal{P}$ a relatively representable moduli problem affine over \EllR. Then conclusions (1), (2), (3) of (\ref{cor:8.4.10}) hold, with $[U, a]$ replaced by $[\mu_4]$.
\end{kcor}

% ------------------------------------------------------------------------------
% Section 8.5
% ------------------------------------------------------------------------------
\section{Base change for coarse moduli schemes}
\label{sec:8.5}

\kmitem\label{item:8.5.1}
Let $R$ be a ring, $\mathcal{P}$ a relatively representable moduli problem which is affine over \EllR, and $U$ an affine open set of $\Spec(R[j])$ on which $j(j-1728)$ is invertible. For any extension of scalars $R \to R'$, we have natural morphisms of $R'[j]$-schemes
\begin{equation}\label{eq:8.5.1.1}
\Mcoarse(\mathcal{P}) \tens{R} R' \to \Mcoarse(\mathcal{P} \tens{R} R') ,
\end{equation}
\begin{equation}\label{eq:8.5.1.2}
\Mcoarse(\mathcal{P}|U) \tens{R} R' \to \Mcoarse((\mathcal{P}|U) \tens{R} R') ,
\end{equation}
the second of which is simply the restriction of the first to the open set $U$.

\kmitem\label{item:8.5.2}
If the first of these maps is an isomorphism for any $R \to R'$, we say that $\mathcal{P}$ satisfies coarse base change. If the second is always an isomorphism, we say that $\mathcal{P}|U$ satisfies coarse base change.

\begin{kthm}\label{thm:8.5.3}
\begin{enumerate}[label=(\Roman*), leftmargin=*]
\item $\mathcal{P}$ satisfies coarse base change if any of the following conditions holds:
\begin{enumerate}[label=(\arabic*), leftmargin=*]
\item $\mathcal{P}$ is representable.
\item The integer $6 = 2 \times 3$ is invertible in $R$.
\item $\mathcal{P}$ is the quotient of a representable $\mathcal{P}'$ by a finite group $G$ whose order is invertible in $R$.
\end{enumerate}
\item $\mathcal{P}|U$ satisfies coarse base change if any of the following conditions holds:
\begin{enumerate}[label=(\arabic*), start=0, leftmargin=*]
\item $\mathcal{P}$ satisfies coarse base change
\item The integer $2$ is invertible in $R$.
\item For any $E/S/R$ which lies in $(\mathrm{Ell}/U/R)$, the automorphism $[-1]_E$ acts without fixed points on the set $\mathcal{P}(E/S)$.
\item $R$ is a Dedekind domain with fraction field of characteristic zero, for every $E/S/R$, the automorphism $[-1]_E$ acts trivially on the set $\mathcal{P}(E/S)$, and $\mathcal{P}$ is flat over $R$.
\end{enumerate}
\end{enumerate}
\end{kthm}

\begin{proof}
Part I has already been proven (\ref{prop:8.1.6}). Let us prove II. If $2$ is invertible, then $\mathcal{P}|U$ is the quotient of a representable problem, namely $(\mathcal{P}, [U, a])$, by $\pm 1$, and the result follows by I(3). Under the hypothesis II(2), $\mathcal{P}|U$ is representable (because rigid). To prove II(3), it suffices to do so after the faithfully flat base change
\[
R \to R[1/2] \oplus ( \bigoplus_{v|2} R_v )
\]
where the sum is over the places $v$ of $R$ of residue characteristic $2$, and $R_v$ denotes the $v$-adic completion. Over $R[1/2]$, II(1) applies. So we are reduced to considering an $R_v$. By (\ref{prop:A7.1.4}), it suffices to show that for every $n \ge 1$,
\[
\mathcal{P} \tens{R} (R_v/\pi_v^n R_v)|U
\]
satisfies coarse base change over the ring $R_v/\pi_v^n R_v$. But $2$ is nilpotent in this ring, so that by (\ref{cor:8.4.12}), the coarse moduli scheme in question is a fine moduli scheme, to which I(1) applies
\[
\begin{gathered}[b]
\Mfrak(\mathcal{P} \tens{R} (R_v/\pi_v^n R_v), [\mu_4]) \\[3pt]
\downarrow \\[3pt]
\Mcoarse(\mathcal{P} \tens{R} (R_v/\pi_v^n R_v))|U .
\end{gathered} \qedhere
\]
\end{proof}

\kmitem\label{item:8.5.4}Summarizing Table

In the table below, $N \ge 1$ is an integer, and $U = U_N$ is any affine open set of the $j$-line over $\Z$ such that
\[
\left\{
\begin{array}{@{}l@{}}
\text{$j(j-1728)$ is invertible on $U_N$.} \\[10pt]
\parbox[t]{0.60\linewidth}{For any prime $p$ dividing $N$, and any supersingular $E/k$ with $k$ a field of characteristic $p$, the $j$-invariant $j(E/k)$ does not lie in $U_N(k)$.}
\end{array}
\right.
\]

\begin{center}
\resizebox{\linewidth}{!}{%
\setlength{\tabcolsep}{3pt}%
\begin{tabular}{@{}c|c|c@{}}
\hline
Problem $\mathcal{P}$ & $\begin{array}{@{}c@{}} \text{Coarse base change} \\ \mathcal{P}|U\ \text{over}\ \Z \end{array}$ & $\begin{array}{@{}c@{}} \text{Justification} \\ \text{over}\ \Z_2 \end{array}$ \\
\hline
$[\Gamma(N)]$ & YES & $\begin{array}{@{}l@{}} \text{if $N=1,2$; II(3)} \\ \text{if $4|N$; II(2)} \\ \text{if odd $p|N$; II(2)} \end{array}$ \\
\cline{1-2}
$[\mathrm{bal.}\ \Gamma_1(N)]$ & $\begin{array}{@{}l@{}} \text{YES \ if $N \neq 4$} \\ \text{? \ if $N = 4$} \end{array}$ & $\begin{array}{@{}l@{}} \text{if $N=1,2$; II(3)} \\ \text{if $8|N$; II(2)} \\ \text{if odd $p|N$; II(2)} \end{array}$ \\
\cline{1-2}
$[\Gamma_1(N)]$ & $\begin{array}{@{}l@{}} \text{YES \ if $N \neq 4,8,16\cdots$} \\ \text{? \ if $N = 4,8,16\cdots$} \end{array}$ & $\begin{array}{@{}l@{}} \text{if $N=1,2$; II(3)} \\ \text{if odd $p|N$; II(2)} \end{array}$ \\
\cline{1-2}
$[\Gamma_0(N)]$ & YES & II(3) \\
\cline{1-2}
$[\Gamma_0(p^n); a, b] \text{ with } a \ge 1, b \ge 1$ & $\begin{array}{@{}l@{}} \text{YES \ if $p$ odd} \\ \text{YES \ if $p=2$, $n \ge 3$, $a \ge 2$, $b \ge 2$} \\ \text{YES \ if $p=2$, $a=b=1$} \\ \text{? \ in other cases} \end{array}$ & $\begin{array}{@{}l@{}} \text{II(2)} \\ \text{II(2)} \\ \text{II(3)} \end{array}$ \\
\hline
\end{tabular}}
\end{center}

\section{Cusps by normalization near infinity; compactified coarse moduli schemes}
\label{sec:8.6}

\kmitem\label{item:8.6.1}
In this section, we will assume that $R$ is a noetherian ring which is regular and \emph{excellent} (e.g., $R$ is excellent if it is of finite type over $\Z$, over a field, over a mixed-characteristic discrete valuation ring \dots\ cf.\ \cite[IV, 7.8]{EGA}).

\kmitem\label{item:8.6.2}
Let $\mathcal{P}$ be a moduli problem on \EllR\ which satisfies the following two axioms:
\begin{enumerate}[label=\textbf{C\arabic*.}]
\item $\mathcal{P}$ is relatively representable and finite over \EllR
\item $\Mcoarse(\mathcal{P})$ is normal near infinity, in the sense that there exists a monic polynomial $f(j) \in R[j]$ such that over the open set $U \subset \Aff^1_R$ where $f$ is invertible, $\Mcoarse(\mathcal{P})|U$ is normal.
\end{enumerate}
Then $\Mcoarse(\mathcal{P})$ is finite over $\Aff^1_R = \Spec(R[j])$, and the open set $\Mcoarse(\mathcal{P})|U$ is normal. Notice that if $\mathcal{P}$ itself is normal near infinity, i.e., if $\mathcal{P}|U$ is normal, then $\Mcoarse(\mathcal{P})$ is normal near infinity.

\kmitem\label{item:8.6.3}
We now extend the $\Aff^1_R$-scheme $\Mcoarse(\mathcal{P})$ to a $\PP^1_R$-scheme $\overline{\Mcoarse}(\mathcal{P})$ by ``normalizing near infinity.'' To be precise, let us temporarily denote by $D$ the divisor $f = 0$ in $\Aff^1_R$. Because $f$ is monic, the open set
\[
\overline{U} \xlongequal{\mathrm{dfn}} \PP^1_R - D
\]
is a neighborhood of the $\infty$-section in $\PP^1_R$. We define
\[
\begin{tikzcd}
\overline{\Mcoarse}(\mathcal{P}) \arrow[d] \\
\PP^1_R
\end{tikzcd}
\]
as follows. Over $\Aff^1_R$, it coincides with $\Mcoarse(\mathcal{P})$. Over $\overline{U}$, it is the \emph{normalization} of $\overline{U}$ in $\Mcoarse(\mathcal{P})|U$:
\begin{equation}\label{eq:8.6.3.1}
\begin{tikzcd}[column sep=normal]
\Mcoarse(\mathcal{P}) \arrow[d] & \Mcoarse(\mathcal{P})|U \arrow[l, hook'] \arrow[r, hook] \arrow[d] & \Mcoarse(\mathcal{P})|\overline{U} \arrow[d] \\
\Aff^1_R & U \arrow[l, hook'] \arrow[r, hook] & \overline{U} .
\end{tikzcd}
\end{equation}
These two pieces agree over $\overline{U} \cap \Aff^1_R = U$, precisely because $\Mcoarse(\mathcal{P})|U$ is already normal, and finite over $U$.

Thus $\overline{\Mcoarse}(\mathcal{P})$ is the unique scheme finite over $\PP^1_R$, which agrees with $\Mcoarse(\mathcal{P})$ over $\Aff^1_R$ and which over some open neighborhood $\overline{U}$ of the $\infty$-section in $\PP^1_R$ is normal (cf.\ \cite[36.6]{Nag}, \cite[IV, 7.8.3, (vi)]{EGA}).

The scheme of cusps of $\mathcal{P}$ is by definition the reduced finite $R$-scheme
\begin{equation}\label{eq:8.6.3.2}
\mathrm{Cusps}(\mathcal{P}) \xlongequal{\mathrm{dfn}} (\overline{\Mcoarse}(\mathcal{P}) - \Mcoarse(\mathcal{P}))^{\mathrm{red}} = (j^{-1}(\infty))^{\mathrm{red}} .
\end{equation}
The formal completion of $\overline{\Mcoarse}(\mathcal{P})$ along its cusps, or equivalently along the divisor defined by the equation ``$1/j = 0$'', is denoted $\widehat{\mathrm{Cusps}}(\mathcal{P})$. We have a cartesian diagram
\begin{equation}\label{eq:8.6.3.3}
\begin{tikzcd}[column sep=small]
\widehat{\mathrm{Cusps}}(\mathcal{P}) \arrow[rr, hook] \arrow[d] & & \overline{\Mcoarse}(\mathcal{P}) \arrow[d] \\
\Spec(R[[1/j]]) \arrow[r, hook] & \Spec(R[1/j]) \arrow[r, hook] & \PP^1_R .
\end{tikzcd}
\end{equation}

\begin{kprop}\label{prop:8.6.4}
Let $R$ be an excellent noetherian regular ring, $\mathcal{P}$ a moduli problem on \EllR\ satisfying C1 and C2, and $G$ a finite group operating on $\mathcal{P}$. Then $\mathcal{P}/G$ satisfies C1 and C2 (by \ref{thm:7.1.3}, (5), (6)), and
\begin{enumerate}[(1)]
\item $G$ acts on $\overline{\Mcoarse}(\mathcal{P})$, covering its trivial action on $\PP^1_R = \overline{\Mcoarse}([\Gamma(1)])$.
\item We have an isomorphism $\overline{\Mcoarse}(\mathcal{P})/G \xrightarrow{\sim} \overline{\Mcoarse}(\mathcal{P}/G)$ of $\PP^1_R$-schemes.
\item We have an isomorphism
\[
\widehat{\mathrm{Cusps}}(\mathcal{P})/G \xrightarrow{\sim} \widehat{\mathrm{Cusps}}(\mathcal{P}/G)
\]
of $R[[1/j]]$-schemes.
\end{enumerate}
\end{kprop}

\begin{proof}
(1) The action of $G$ on $\Mcoarse(\mathcal{P})/\Mcoarse([\Gamma(1)])$ extends by normalization. (2) The isomorphism $\Mcoarse(\mathcal{P})/G \xrightarrow{\sim} \Mcoarse(\mathcal{P}/G)$ of $\Mcoarse([\Gamma(1)])$-schemes extends by normalization (the scheme $(\Mcoarse(\mathcal{P})/G)|\overline{U}$ is normal, finite over $\overline{U}$, and agrees with $\overline{\Mcoarse}(\mathcal{P}/G)$ over $U$). (3) follows from (2) and the above cartesian diagram, which shows that $\widehat{\mathrm{Cusps}}(\mathcal{P})/R[[1/j]]$ is obtained from $\overline{\Mcoarse}(\mathcal{P})/\PP^1_R$ by a flat base extension. Because formation of quotients by finite groups commutes with flat base extension, we get (3).
\end{proof}

\begin{definition}\label{def:8.6.5}
We say that a moduli problem $\mathcal{P}$ on \EllR\ which satisfies C1 and C2 is \emph{smooth near infinity} if it satisfies \emph{both} of the following conditions:
\begin{enumerate}[label=\textbf{S-$\infty$-\arabic*.}]
\item There exists an open neighborhood $V$ of the cusps in $\overline{\Mcoarse}(\mathcal{P})$,
\[
\mathrm{Cusps}(\mathcal{P}) \subset V \subset \overline{\Mcoarse}(\mathcal{P})
\]
which is smooth over $R$ (necessarily of relative dimension one).
\item The scheme of cusps $\mathrm{Cusps}(\mathcal{P})$ is finite etale over $R$.
\end{enumerate}
In practice, it is convenient to replace S-$\infty$-1 by the equivalent
\begin{enumerate}[label=\textbf{S-$\infty$-1 bis.}]
\item The formal completion $\widehat{\mathrm{Cusps}}(\mathcal{P})$ of $\overline{\Mcoarse}(\mathcal{P})$ along its cusps is formally smooth over $R$.
\end{enumerate}
\end{definition}

\begin{kprop}\label{prop:8.6.6}
Let $R$ be an excellent noetherian regular ring, and $\mathcal{P}$ a moduli problem on \EllR\ which is relatively representable. Suppose that:
\begin{enumerate}[(1)]
\item $\mathcal{P}$ is finite over \EllR, and $\Mcoarse(\mathcal{P})$ is normal near infinity.
\item $\mathcal{P}$ satisfies coarse base change near infinity, in the sense that there exists a monic polynomial $f(j) \in R[j]$, such that for $U$ the open set where $j(j-1728)f(j)$ is invertible, $\mathcal{P}|U$ satisfies coarse base change. $\langle$Recall that this is automatic if $\mathcal{P}$ is representable, or if $2$ is invertible, or \dots\ (cf.\ \ref{thm:8.5.3})$\rangle$.
\item $\mathcal{P}$ is smooth near infinity.
\end{enumerate}
Then for any extension of scalars $R \to R'$ with $R'$ an excellent noetherian regular ring, $\mathcal{P} \tens{R} R'$ also satisfies (1), (2), (3), and we have
\begin{enumerate}[(a)]
\item for some open neighborhood $\overline{U}$ of the $\infty$-section of $\PP^1_R$,
\[
\overline{\Mcoarse}(\mathcal{P}) \tens{R} R' | \overline{U} \xrightarrow{\sim} \overline{\Mcoarse}(\mathcal{P} \tens{R} R') | \overline{U} .
\]
\nopagebreak
\item $\widehat{\mathrm{Cusps}}(\mathcal{P}) \tens{R} R' \xrightarrow{\sim} \widehat{\mathrm{Cusps}}(\mathcal{P} \tens{R} R')$.
\nopagebreak
\item $\mathrm{Cusps}(\mathcal{P}) \tens{R} R' \xrightarrow{\sim} \mathrm{Cusps}(\mathcal{P} \tens{R} R')$.
\end{enumerate}
\end{kprop}

\begin{proof}
Enlarging the monic polynomial $f$, we may assume that the open neighborhood $V$ of the cusps in $\overline{\Mcoarse}(\mathcal{P})$ is the inverse image of $\overline{U} = \PP^1_R - (\text{the divisor } j(j-1728)f(j))$. Then $V \cap \Mcoarse(\mathcal{P})$ is $\Mcoarse(\mathcal{P})|U$, which by hypothesis is smooth over $R$. Therefore $(\Mcoarse(\mathcal{P})|U) \tens{R} R'$ is smooth over $R'$. Because $\mathcal{P}|U$ satisfies coarse base change, this says that $\Mcoarse(\mathcal{P} \tens{R} R')|U$ is smooth over $R'$, hence normal. This already shows that $\mathcal{P} \tens{R} R'$ satisfies (1) and (2). The scheme $\overline{\Mcoarse}(\mathcal{P}) \tens{R} R' | \overline{U}$ is smooth over $R'$, hence normal. As it is finite over $\overline{U} \tens{R} R'$, and coincides over $U \tens{R} R'$ with $\Mcoarse(\mathcal{P}) \tens{R} R' | U = \Mcoarse(\mathcal{P} \tens{R} R') | U$, it must be $\overline{\Mcoarse}(\mathcal{P} \tens{R} R') | \overline{U}$. This proves that $\mathcal{P} \tens{R} R'$ is smooth near infinity, and that (a) holds. If the cusps of $\mathcal{P}$ are finite etale over $R$, they are a Cartier divisor $Z$ in the smooth $R$-curve $\overline{\Mcoarse}(\mathcal{P})|\overline{U}$. Then $Z \tens{R} R'$ is finite etale over $R'$, it is reduced, and so it is the cusps of $\mathcal{P} \tens{R} R'$ inside the smooth $R'$-curve $(\overline{\Mcoarse}(\mathcal{P})|\overline{U}) \tens{R} R' = \overline{\Mcoarse}(\mathcal{P}) \tens{R} R' | \overline{U}$. Therefore (b) and (c) hold.
\end{proof}

\begin{kprop}\label{prop:8.6.7}
Let $R$ be an excellent noetherian regular ring, and let $R'$ be an excellent noetherian regular $R$-algebra which is flat over $R$, and such that the fibers of $\Spec(R') \to \Spec(R)$ are all geometrically regular (e.g., $R'$ etale over $R$, or smooth over $R$, or a completion of $R$ at a closed point of $\Spec(R)$, \dots). Let $\mathcal{P}$ be a moduli problem on \EllR\ satisfying C1 and C2, i.e., $\mathcal{P}$ is relatively representable and finite over \EllR, and $\Mcoarse(\mathcal{P})$ is normal near infinity. Then
\begin{enumerate}[(1)]
\item The moduli problem $\mathcal{P} \tens{R} R'$ on $(\mathrm{Ell}/R')$ satisfies C1 and C2.
\item The natural extension of scalars map of coarse moduli schemes defines an isomorphism of $R'$-schemes
\[
\Mcoarse(\mathcal{P} \tens{R} R') \xrightarrow{\sim} \Mcoarse(\mathcal{P}) \tens{R} R' .
\]
\item The above isomorphism extends to define an isomorphism of compactified coarse moduli schemes
\[
\overline{\Mcoarse}(\mathcal{P} \tens{R} R') \xrightarrow{\sim} \overline{\Mcoarse}(\mathcal{P}) \tens{R} R' .
\]
\end{enumerate}
\end{kprop}

\begin{proof}
The map in (2) is an isomorphism because $R \to R'$ is flat. Because of the ``regular fibers'' hypothesis, $\Mcoarse(\mathcal{P}) \tens{R} R'$ is still normal near infinity, so $\mathcal{P} \tens{R} R'$ satisfies C1 and C2. The isomorphism (3) is the fact that both sides are $R'$-schemes finite over the projective $j$-line, normal near infinity, and which agree over the affine $j$-line.
\end{proof}

\begin{kthm}\label{thm:8.6.8}
Let $R$ be an excellent noetherian regular domain, whose fraction field has characteristic zero. Let $\mathcal{P}$ be a representable moduli problem on \EllR\ which is finite over \EllR. Suppose that for some open set $U \subset \Spec(R[j])$ defined by inverting a monic polynomial in $j$, $\mathcal{P}|U$ is finite etale over $[U]$. Then
\begin{enumerate}[(1)]
\item over the open set $U' = U[1/(j(j-1728))]$, the map
\[
\Mfrak(\mathcal{P}|U') \xrightarrow{\ j\ } U'
\]
is finite etale, and tamely ramified along the infinity section of the projective $j$-line.
\item The compactified moduli scheme $\overline{\Mfrak}(\mathcal{P})$ is smooth over $R$ of relative dimension one in a neighborhood of the cusps, the scheme of cusps is finite etale over $R$, and in a neighborhood of the cusps, the invertible sheaf
\[
\Omega^1_{\overline{\Mfrak}(\mathcal{P})/R}(\log \mathrm{cusps})
\]
of one-forms with at worst simple poles along the cusps is free on the one-form $\mathrm{dlog}(1/j)$.
\item For any excellent noetherian regular $R$-algebra $R'$, statements (1) and (2) remain true for $\mathcal{P} \tens{R} R'$, and we have isomorphisms of $\PP^1_R$-schemes
\[
\overline{\Mfrak}(\mathcal{P}) \tens{R} R' \xrightarrow{\sim} \overline{\Mfrak}(\mathcal{P} \tens{R} R') .
\]
\end{enumerate}
\end{kthm}

\begin{proof}
That $\Mfrak(\mathcal{P}|U') \to U'$ is finite etale has already been proven (\ref{cor:8.4.5}). That it is tamely ramified along the infinity section results from the fact that $R$ has generic characteristic zero. Once we know $\Mfrak(\mathcal{P}|U') \to U'$ is finite etale and tame along infinity, (2) and (3) result by Abhyankar's lemma (cf.\ \cite[I, Exp.~XIII, 5.5]{SGA}).
\end{proof}

% ------------------------------------------------------------------------------
% Section 8.7
% ------------------------------------------------------------------------------
\section{Interlude: The groups $\mathbf{T}[N]$ and $\mathbf{T}$}
\label{sec:8.7}

\kmitem\label{item:8.7.1}
Over the ring $\Z[q,q^{-1}]$ of Laurent polynomials over $\Z$ in one variable, we will define for every integer $N \ge 1$ a finite flat commutative group-scheme $\mathbf{T}[N]$ of rank $N^2$ which is killed by $N$. As a scheme, $\mathbf{T}[N]$ is the disjoint union of $N$ schemes $\mathbf{T}_0[N], \dots, \mathbf{T}_{N-1}[N]$, where
\begin{equation}\label{eq:8.7.1.1}
\mathbf{T}_i[N] = \Spec(\Z[q,q^{-1}][X]/(X^N - q^i)) .
\end{equation}
For any $\Z[q,q^{-1}]$-algebra $R$ with connected spectrum, we have
\begin{equation}\label{eq:8.7.1.2}
\begin{gathered}
\mathbf{T}[N](R) = \text{pairs } (X, i/N) \text{ with } 0 \le i \le N-1 \text{ and with} \\
X \in R \text{ satisfying } X^N = q^i .
\end{gathered}
\end{equation}
The group structure is given by defining
\begin{equation}\label{eq:8.7.1.3}
(X, i/N) \cdot (Y, j/N) = 
\begin{cases}
(XY, (i+j)/N) & \text{if } i+j \le N-1 \\
(XY/q, (i+j-N)/N) & \text{if } i+j \ge N .
\end{cases}
\end{equation}
The elements of the form $(X,0)$ form a subgroup $\mu_N(R)$, and we have a short exact sequence
\begin{equation}\label{eq:8.7.1.4}
\begin{gathered}
0 \longrightarrow \mu_N \xrightarrow{a_N} \mathbf{T}[N] \xrightarrow{b_N} \Z/N\Z \longrightarrow 0 \\
\zeta \longmapsto (\zeta, 0) \\
(X, i/N) \longmapsto i \bmod N .
\end{gathered}
\end{equation}
This exact sequence splits over any $R$ containing an $N$'th root of $q$, say $q^{1/N}$, for then we may write
\begin{equation}\label{eq:8.7.1.5}
(X, i/N) = (Xq^{-i/N}, 0) \times (q^{i/N}, i/N) .
\end{equation}
Thus over $\Z[q,q^{-1}][q^{1/N}]$, we have
\begin{equation}\label{eq:8.7.1.6}
\mathbf{T}[N] \simeq \mu_N \times \Z/N\Z .
\end{equation}
There is a unique alternating\footnote{i.e., for any $\Z[q,q^{-1}]$-algebra $R$, and any $P \in \mathbf{T}[N](R)$, we have $e_n(P,P) = 1$.} pairing of $\Z[q,q^{-1}]$-group-schemes
\begin{equation}\label{eq:8.7.1.7}
e_N : \mathbf{T}[N] \times \mathbf{T}[N] \to \mu_N ,
\end{equation}
which is compatible with the above extension structure in the following sense: for every $\Z[q,q^{-1}]$-algebra $R$, every $\zeta \in \mu_N(R)$ and every $P \in \mathbf{T}[N](R)$, we have
\begin{equation}\label{eq:8.7.1.8}
e_N(a_N(\zeta), P) = \zeta^{b_N(P)} .
\end{equation}
(It is unique because these formulas and the requirement it be alternating uniquely determine it after the f.p.p.f. base extension to $\Z[q,q^{-1}][q^{1/N}]$. The pairing is given explicitly by the formula
\begin{equation}\label{eq:8.7.1.9}
e_N((X, i/N), (Y, j/N)) = X^j / Y^i .
\end{equation}
We will refer to the exact sequence
\begin{equation}\label{eq:8.7.1.10}
0 \longrightarrow \mu_N \xrightarrow{a_N} \mathbf{T}[N] \xrightarrow{b_N} \Z/N \longrightarrow 0
\end{equation}
as the canonical extension-structure on $\mathbf{T}[N]$, and to the alternating pairing
\[
e_N : \mathbf{T}[N] \times \mathbf{T}[N] \to \mu_N
\]
as the $e_N$-pairing.

\kmitem\label{item:8.7.2}
We will now define a smooth one-dimensional commutative group-scheme $\mathbf{T}$ over $\Z[q,q^{-1}]$, locally of finite type, which sits in an f.p.p.f. short exact sequence of group-schemes over $\Z[q,q^{-1}]$
\begin{equation}\label{eq:8.7.2.1}
0 \to \mathbb{G}_m \to \mathbf{T} \to \Q/\Z \to 0
\end{equation}
in such a way that $\mathbf{T}[N]$ is the kernel of $N$ in $\mathbf{T}$, for every integer $N \ge 1$. As a scheme, $\mathbf{T}$ is the disjoint union of schemes $\mathbf{T}_a$, indexed by $a$'s running over all rational numbers in the interval $[0,1)$, where each $\mathbf{T}_a$ is just the scheme $\mathbb{G}_m$. Thus over any $\Z[q,q^{-1}]$-algebra $R$ with connected spectrum, we have
\begin{equation}\label{eq:8.7.2.2}
\mathbf{T}(R) = \text{pairs } (X, \alpha) \text{ with } X \in R^{\times} \text{ and } \alpha \in \Q \cap [0,1) .
\end{equation}
The group-structure is defined by
\begin{equation}\label{eq:8.7.2.3}
(X, \alpha) \cdot (Y, \beta) = 
\begin{cases}
(XY, \alpha+\beta) & \text{if } \alpha+\beta < 1 \\
\left(\dfrac{XY}{q}, \alpha+\beta-1\right) & \text{if } \alpha+\beta \ge 1 .
\end{cases}
\end{equation}
Let us temporarily denote by $[a]$ the integral part of the rational number $\alpha$, and by $\langle \alpha \rangle$ its fractional part:
\begin{equation}\label{eq:8.7.2.4}
\begin{cases}
\alpha = [\alpha] + \langle \alpha \rangle & \text{with} \\
[\alpha] \in \Z, \ 0 \le \langle \alpha \rangle < 1 .
\end{cases}
\end{equation}
For any $\Z[q,q^{-1}]$-algebra $R$ with connected spectrum, we have a fundamental short exact sequence of groups
\stepcounter{equation}
\begin{equation}\label{eq:8.7.2.5}
\tag*{(\theequation)(*)}
0 \longrightarrow \Z \xrightarrow{a} \mathbb{G}_m(R) \times \Q \xrightarrow{\pi} \mathbf{T}(R) \longrightarrow 0
\end{equation}
defined by
\begin{equation}\label{eq:8.7.2.6}
\begin{gathered}
a(n) = (q^n, n) \\
\pi(X, \alpha) = (X/q^{[\alpha]}, \langle \alpha \rangle) .
\end{gathered}
\end{equation}
Thus (\ref{eq:8.7.2.1}) is the ``pushout'' of the short exact sequence
\begin{equation}\label{eq:8.7.2.7}
0 \to \Z \to \Q \to \Q/\Z \to 0
\end{equation}
of constant group-schemes over $\Z[q,q^{-1}]$ by the tautological homomorphism of group-schemes over $\Z[q,q^{-1}]$
\[
\Z \to \mathbb{G}_m ;
\]
\[
n \mapsto q^{-n} \in \mathbb{G}_m(\Z[q,q^{-1}]) .
\]
The short exact sequence $(*)$ allows for easy computation. Thus given a point $(X, \alpha) \in \mathbf{T}(R)$, i.e., $X \in R^{\times}$, $\alpha \in \Q \cap [0,1)$, we have, for every integer $N$, the formula
\begin{equation}\label{eq:8.7.2.8}
N \cdot (X, \alpha) = (X^N/q^{[N\alpha]}, \langle N\alpha \rangle) .
\end{equation}

\kmitem\label{item:8.7.3}
In particular, $(X, \alpha)$ is killed by $N$ if and only if $\alpha = i/N$ for some $0 \le i \le N-1$ and $X^N = q^i$. Thus as promised we find $\mathbf{T}[N] = \mathbf{T}[N]$ compatibly with the canonical extension structures on $\mathbf{T}$ and on $\mathbf{T}[N]$.

\kmitem\label{item:8.7.4}
Over any $\Z[q,q^{-1}]$-algebra $R$ containing a compatible system of $N$'th roots $q^{1/N}$ of $q$ for every $N$, (i.e., for every integer $N \ge 1$, we are given $Y_N \in R^{\times}$ such that $Y_1 = q$ and $(Y_{NM})^M = Y_N$ for every $M,N \ge 1$), the inverse image $\mathbf{T}_R$ of $\mathbf{T}$ on $\Spec(R)$ has a product decomposition as $R$-group-scheme
\begin{equation}\label{eq:8.7.4.1}
\mathbf{T}_R \xrightarrow{\sim} \mathbb{G}_m \times \Q/\Z .
\end{equation}
defined on $R'$-valued points, for any $R$-algebra $R'$ with connected spectrum, by
\addtocounter{equation}{-1}
\begin{equation}
\begin{tikzcd}
(X, a) \arrow[r, mapsto] & (X/q^a, a) \\
\mathbf{T}(R') \arrow[u, "\rotatebox{90}{$\in$}" description, phantom] & \mathbb{G}_m(R') \times \Q/\Z \arrow[u, "\rotatebox{90}{$\in$}" description, phantom]
\end{tikzcd}
\end{equation}
where ``$q^a$'' means that we write $a = i/N$ with $0 \le i \le N-1$ and then define $q^a$ to be $(q^{1/N})^i$. Because the $q^{1/N}$'s are compatible, this $q^a$ is well defined independently of the representation of $a$ as a fraction.

\begin{kthm}\label{thm:8.7.5}
There exists a faithfully flat $\Z[q,q^{-1}]$-algebra $R$, an elliptic curve $E/R$, and an isomorphism of ind-group-schemes over $R$
\[
\mathbf{T}_{\mathrm{torsion}} \tens{\Z[q,q^{-1}]} R \xrightarrow{\sim} E_{\mathrm{tors}} ,
\]
such that for every $N \ge 1$, the isomorphism on $N$-division points
\[
\mathbf{T}[N] \otimes R \xrightarrow{\ \sim\ } E[N]
\]
is compatible with $e_N$-pairings.
\end{kthm}

\begin{proof}
Over the faithfully flat $\Z[q,q^{-1}]$-algebra
\[
\varinjlim_{N \ge 1} \Z[q,q^{-1}][X_N]/((X_N)^{N!} - q)
\]
with transition maps $X_N \mapsto (X_{N+1})^{N+1}$, we have a compatible system of $N$'th roots of $q$, so $\mathbf{T}$ becomes isomorphic to the product $\mathbb{G}_m \times (\Q/\Z)$, and $\mathbf{T}_{\mathrm{tors}}$ therefore becomes isomorphic to the product
\[
\mu_{\infty} \times \Q/\Z .
\]
As this group with all its $e_N$-pairings begins life over $\Z$, we are reduced to showing
\renewcommand{\qedsymbol}{}
\end{proof}

\begin{klem}\label{lem:8.7.6}
There exists a faithfully flat $\Z$-algebra $R$, an elliptic curve $E/R$, and an isomorphism of ind-group-schemes over $R$
\[
\phi_{\infty} : \mu_{\infty} \times \Q/\Z \xrightarrow{\sim} E_{\mathrm{tors}}
\]
which is compatible with all $e_N$-pairings.
\end{klem}

\begin{proof}
Suppose we have constructed an $E/R$ and an isomorphism $\phi_{\infty}$. Then we can always modify $\phi_{\infty}$ by an automorphism of $\mu_{\infty} \times \Q/\Z$ of the form (auto.\ of $\mu_{\infty}$) $\times$ (id) to render it compatible with the $e_N$-pairings. For let $e'_N$ be the pairings on the $\mu_N \times \Z/N\Z$'s obtained via $\phi_{\infty}$ from the $e_N$-pairings on the $E[N]$'s. Because they are alternating, we necessarily have
\[
\begin{aligned}
& e'_N((\zeta, i), (\eta, j)) \\
& \qquad = e'_N((\zeta^j/\eta^i, 0), (1,1)) .
\end{aligned}
\]
\emergencystretch=6em
[Because $R$ is flat over $\Z$, it suffices to verify this over $R[1/N]$, where it is physically obvious that the ``other terms'' $e'_N((\zeta,0), (\eta,0))$ and $e'_N((1,i), (1,j))$ are both equal to $1$, because $e'_N$ is alternating in the strong sense $e'_N(x,x) = 1$.] Because $e'_N$ must define an auto-duality on $\mu_N \times \Z/N\Z$ (this being true for $e_N$ on $E[N]$), the map
\[
\begin{gathered}
\mu_N \longrightarrow \mu_N \\
\zeta \longmapsto e'_N((\zeta,0), (1,1))
\end{gathered}
\]
\emergencystretch=2em
must be an automorphism $a_N$ of $\mu_N$. For variable $N$, the compatibilities satisfied by the $e_N$-pairings on the $E[N]$ give compatibilities on the $e'_N$-pairings which force these automorphisms $a_N$ to fit together and provide a single automorphism $a_{\infty}$ of $\mu_{\infty}$. Then if we compose $\phi_{\infty}$ with ($a_{\infty}^{-1}$ on $\mu_{\infty}$) $\times$ (id), we get a ``new'' $\phi_{\infty}$ which is compatible with $e_N$-pairings.

We now turn to the construction of $(E/R, \phi_{\infty})$ with $R$ faithfully flat over $\Z$. For each prime number $p$, choose an algebraically closed field $k(p)$ of characteristic $p$. Denote by $W(k(p))$ its ring of Witt vectors.

Then the product ring
\[
\prod_p W(k(p))
\]
is faithfully flat over $\Z$ (since it's torsion free, and non-zero mod $p$ for every $p$). Thus it suffices to construct, separately for each $p$, an elliptic curve $E/W(k(p))$ together with a $\phi_{\infty}$.

Choose an ordinary elliptic curve $E_0$ over $k(p)$. Because $k(p)$ is algebraically closed, and $E_0$ is ordinary, there exists an isomorphism of $p$-divisible groups over $k(p)$
\[
\phi_{0,p} : E_0[p^{\infty}] \simeq \mu_{p^{\infty}} \times \Q_p/\Z_p .
\]
The prime-to-$p$ torsion subgroup is etale, so again because $k(p)$ is algebraically closed, there exists an isomorphism over $k(p)$
\[
\phi_{0,\, \text{not } p} : E_0[\text{prime-to-}p \text{ torsion}] \simeq \mu_{\text{prime to } p} \times \Bigl(\prod_{\ell \neq p} \Q_{\ell}/\Z_{\ell}\Bigr) .
\]
These isomorphisms together form an isomorphism over $k(p)$
\[
\phi_{0,\infty} : (E_0)_{\mathrm{torsion}} \xrightarrow{\sim} \mu_{\infty} \times \Q/\Z .
\]
Let $E/W(k(p))$ be the \emph{canonical lifting} \`a la Serre--Tate of $E_0/k(p)$. Then by its very definition, any isomorphism $\phi_{0,p}$ lifts to an isomorphism of $p$-divisible groups over $W(k(p))$:
\[
\phi_p : E[p^{\infty}] \simeq \mu_{p^{\infty}} \times \Q_p/\Z_p .
\]
The isomorphism $\phi_{0, \text{not } p}$ also lifts uniquely, and the two together define the inverse of an isomorphism
\[
\mu_{\infty} \times \Q/\Z \xrightarrow{\sim} E_{\mathrm{torsion}} . \qedhere
\]
\end{proof}

% ------------------------------------------------------------------------------
% Section 8.8
% ------------------------------------------------------------------------------
\section{Relation to the Tate curve}
\label{sec:8.8}

For any ring $R$, we denote by $R((q))$ the ring of finite-tailed Laurent series over $R$ in one variable $q$. One knows (\cite{Roq}, \cite{DeRa}, \cite{K2}) that over the ring $\Z((q))$ there is an elliptic curve $\Tate(q)$ together with a nowhere vanishing invariant one-form $\wcan$, which satisfies all of the following properties T.1 through T.4.

\tateprop\label{tate:1}
{\renewcommand{\arraystretch}{1.35}%
\(\displaystyle
\left\{
\begin{array}{l}
\displaystyle j(\Tate(q)) = \frac{1}{q} + 744 + \cdots \\[2pt] \hline
\displaystyle c_4(\Tate(q), \wcan) = 1 + 240 \sum_{n \ge 1} q^n \sum_{d | n} d^3 \\[2pt] \hline
\displaystyle c_6(\Tate(q), \wcan) = 1 - 504 \sum_{n \ge 1} q^n \sum_{d | n} d^5 \\[2pt] \hline
\displaystyle \Delta(\Tate(q), \wcan) = q \cdot \prod_{n \ge 1} (1 - q^n)^{24} .
\end{array}
\right.
\)\par}

\tateprop\label{tate:2}
There exists a unique isomorphism $\phi_{\mathrm{can}}$ of formal Lie groups over $\Z((q))$
\[
\phi_{\mathrm{can}} : \widehat{\Tate(q)} \xrightarrow{\ \sim\ } \widehat{\mathbb{G}}_m = \hat{\mathbf{T}}
\]
under which the ``standard'' invariant differential $\mathrm{d}X/X$ on $\mathbb{G}_m \subset \mathbf{T}$ pulls back to $\wcan$
\[
\phi_{\mathrm{can}}^*(\mathrm{d}X/X) = \wcan .
\]

\tateprop\label{tate:3}
For every $N \ge 1$, there exists a unique short exact sequence (as f.p.p.f. abelian sheaves) of finite flat commutative group-schemes over $\Z((q))$
\[
0 \longrightarrow \mu_N \xrightarrow{a_N} \Tate(q)[N] \xrightarrow{b_N} \Z/N\Z \longrightarrow 0 ,
\]
with the following properties:

T.3.a$_N$: for any $\Z((q))$-algebra $R$ in which $N$ is nilpotent, and any $\zeta \in \mu_N(R)$, we have
\[
\phi_{\mathrm{can}}(a_N(\zeta)) = \zeta .
\]
The formula makes sense, because when $N$ is nilpotent, $a_N$ necessarily maps $\mu_N$ into the formal group of $\Tate(q)$.

T.3.b$_N$: for any $\Z((q))$-algebra $R$, any element $\zeta \in \mu_N(R)$, and any element $X \in \Tate(q)[N](R)$, we have the formula
\[
e_N(a_N(\zeta), X) = \zeta^{b_N(X)} .
\]

\tateprop\label{tate:4}
There exists a unique isomorphism of ind-group-schemes on $\Z((q))$
\[
\mathbf{T}_{\mathrm{torsion}} \tens{\Z[q,q^{-1}]} \Z((q)) \xrightarrow{\ \sim\ } \Tate(q)_{\mathrm{torsion}}
\]
which is compatible with the canonical extension structures:
\[
\begin{tikzcd}[column sep=1.7em, row sep=2.6em]
0 \arrow[r] & \mu_{\infty} \arrow[r] \arrow[d, equal] & \mathbf{T}_{\mathrm{torsion}} \arrow[r] \arrow[d, "\cong"] & \Q/\Z \arrow[r] \arrow[d, equal] & 0 \\
0 \arrow[r] & \mu_{\infty} \arrow[r] & \Tate(q)_{\mathrm{torsion}} \arrow[r] & \Q/\Z \arrow[r] & 0 .
\end{tikzcd}
\]
[This isomorphism is unique because it is a priori indeterminate up to a $\Z((q))$-homomorphism from $\Q/\Z$ to $\mu_{\infty}$, i.e., up to an element of the one-element group
\[
\varinjlim_N \mu_N(\Z((q))) \xrightarrow{\ \sim\ } \varinjlim_N \mu_N(\Z) = \{1\} .
\]
It is automatically compatible with $e_N$-pairings, because these pairings are uniquely determined by the fact that they are alternating, and compatible with the extension-structures in the sense of (\ref{eq:8.7.1.8}).]

% ------------------------------------------------------------------------------
% Section 8.9
% ------------------------------------------------------------------------------
\section{Relation with ordinary elliptic curves via the Serre-Tate parameter}
\label{sec:8.9}

Let $k$ be an algebraically closed field of characteristic $p > 0$, $W = W(k)$ its ring of Witt-vectors, $R$ a complete noetherian local $W(k)$-algebra with residue field $k$, $E/R$ an elliptic curve, and $E_0/k$ its special fiber.

Because $k$ is algebraically closed, and $E_0/k$ is ordinary, there exists an isomorphism of ind-group-schemes over $k$
\[
\phi_0 : \mu_{\infty} \times \Q/\Z \xrightarrow{\ \sim\ } (E_0)_{\mathrm{torsion}} ,
\]
which is compatible with all the $e_N$-pairings. Associated to the data $(E/R, \phi_0)$ is the ``Serre-Tate parameter''
\[
q(E/R, \phi_0) \in 1 + \max(R) \subset R^{\times} .
\]
By means of the Serre-Tate parameter, we may view $R$ as a $\Z[q,q^{-1}]$-algebra:
\[
q \mapsto q(E/R, \phi_0) .
\]
According to the Serre-Tate theorem (cf. \cite{K5}), there is a \emph{unique} isomorphism of ind-group-schemes over $R$
\[
\phi : \mathbf{T}_{\mathrm{tors}} \tens{\Z[q,q^{-1}]} R \xrightarrow{\ \sim\ } E_{\mathrm{tors}}
\]
which \emph{lift} the given
\[
\phi_0 : \mu_{\infty} \times \Q/\Z \xrightarrow{\ \sim\ } (E_0)_{\mathrm{tors}}
\]
in the following sense: via the isomorphism of $k$-group-schemes
\[
\mathbf{T} \tens{\Z[q,q^{-1}]} k \xrightarrow{\ \sim\ } \mathbb{G}_m \times \Q/\Z
\]
defined by using the canonical system $\{Y_N = 1\}$ of $N$'th roots of the image of $q$ in $k$, we have a commutative diagram of isomorphisms
\[
\begin{tikzcd}[column sep=5.2em, row sep=1.9em]
\mathbf{T}_{\mathrm{tors}} \otimes k
  \arrow[dr, "\phi \bmod \max(R)", "\sim"]
  \arrow[dd, "\rotatebox{90}{$\sim$}"'] & \\
& (E_0)_{\mathrm{tors}} \\
\mu_{\infty} \times \Q/\Z
  \arrow[ur, "\phi_0", "\sim"'] &
\end{tikzcd}
\]
Because $\phi_0$ is compatible with $e_N$-pairings, its lifting $\phi$ is also compatible [by reduction to the universal case $R = W[[q-1]]$ which is flat over $\Z$, one sees as in the proof of (\ref{lem:8.7.6}) that the possible error appears through a section ``det'' over $R$ of the constant group $\varinjlim (\Z/N\Z)^{\times}$. Since we have ``det $= 1$'' for $\phi_0$, and $\Spec(R)$ is connected, we have ``det $= 1$'' for $\phi$ also].

% Section 8.10
\section{Other universality properties of the groups $\mathbf{T}[N]$}
\label{sec:8.10}

\kmitem\label{item:8.10.1}
Let $S$ be an arbitrary scheme, $N \ge 1$ an integer, and $G/S$ a finite locally free commutative group-scheme over $S$ of rank $N^2$ which is killed by $N$. A bilinear pairing
\[
(\ ,\ ) : G \times G \to \mu_N
\]
(as $S$-group-schemes) is said to be \emph{strongly alternating} if the following two conditions hold

\saprop{1}
for any $S$-scheme $T$, and any point $P \in G(T)$, we have
\[
(P,P) = 1 \in \mu_N(T) .
\]

\saprop{2}
for any $S$-scheme $T$, any finite, locally free of rank $N$ $T$-subgroup-scheme $K \subset G \times T$, and any pair of points $P_1, P_2$ both in $K(T)$, we have
\[
(P_1,P_2) = 1 \in \mu_N(T) .
\]

Recall \eqref{eq:2.8.7.1} that the $e_N$-pairing on $E[N]$, for $E/S$ an elliptic curve, is strongly alternating.

\kmitem\label{item:8.10.2}
Clearly if $(\ ,\ )$ is strongly alternating, then for any short exact (as f.p.p.f. sheaves) sequence of $S$-group-schemes
\[
0 \to K \to G \xrightarrow{\ \pi\ } K' \to 0
\]
with $K$ and $K'$ both locally free of rank $N$, there is an induced bilinear pairing of $S$-group-schemes
\[
\langle\ ,\ \rangle : K \times K' \to \mu_N
\]
characterized by the following property:

for any $S$-scheme $T$, for any points $P \in K(T)$, $Q \in G(T)$, we have
\[
(P,Q) = \langle P, \pi Q \rangle .
\]

\kmitem\label{item:8.10.3}
If $G$ is given with a structure of extension of $S$-groups
\[
0 \to \mu_N \to G \to \Z/N\Z \to 0 ,
\]
we say that a strongly alternating bilinear form $(\ ,\ )$ on $G$ is \emph{compatible} with the extension structure if the induced pairing
\[
\langle\ ,\ \rangle : \mu_N \times \Z/N\Z \to \mu_N
\]
is given by
\[
\langle \zeta, a \rangle = \zeta^a .
\]

\begin{klem}\label{lem:8.10.4}
Let $G/S$ be a finite locally free commutative group-scheme of rank $N^2$ which is killed by $N$, given with an extension structure
\[
0 \to \mu_N \to G \to \Z/N\Z \to 0 .
\]
Then there exists a unique strongly alternating $(\ ,\ )$ on $G$ which is compatible with the extension structure.
\end{klem}

\begin{proof}
Locally f.p.p.f. on $S$, the extension is splittable; then $G \simeq \mu_N \times \Z/N\Z$ and in this case the existence and uniqueness of $(\ ,\ )$ are obvious (cf.~\eqref{eq:8.7.1.9}); by f.p.p.f. descent, this unique $(\ ,\ )$ descends to a unique $(\ ,\ )$ on $G$ itself.
\end{proof}

\begin{kprop}\label{prop:8.10.5}
Let $G/S$ be a finite locally free commutative group-scheme of rank $N^2$ which is killed by $N$, given with an extension structure
\[
0 \to \mu_N \to G \xrightarrow{\ \pi\ } \Z/N\Z \to 0 .
\]
Let $(\ ,\ )$ denote the unique strongly alternating bilinear form on $G$ compatible with this extension structure. Then Zariski locally on $S$, there exists a (non-unique) unit $q \in \mathbb{G}_m(S) = \Hom_{\Sch}(S, \Spec(\Z[q,q^{-1}]))$ and a (non-unique) isomorphism of $S$-groups
\[
\mathbf{T}[N] \tens{\Z[q,q^{-1}]} S \xrightarrow{\ \sim\ } G
\]
which is compatible with extension structures. Any such isomorphism necessarily carries the $e_N$-pairing on $\mathbf{T}[N]$ to the pairing $(\ ,\ )$ on $G$.
\end{kprop}

\begin{proof}
Locally f.p.p.f. on $S$, the exact sequence
\[
0 \to \mu_N \to G \xrightarrow{\ \pi\ } \Z/N\Z \to 0
\]
splits (because $N$ kills $G$, a group-theoretic splitting amounts to specifying a section $g$ of $G$ with $\pi(g) = 1$), so our $G$ is an f.p.p.f. \emph{form} of the product group-scheme $\mu_N \times \Z/N\Z$ with its extension structure. The group $\mathbf{T}[N]$ with its canonical extension structure is \emph{another} form.

The set of $S$-isomorphism classes of forms of $\mu_N \times \Z/N\Z$ with its extension structure is the cohomology group $H^1_{\text{f.p.p.f.}}(S, \mathit{Aut})$, where $\mathit{Aut}$ denotes the $S$-group-scheme whose $T$-valued points are
\[
\resizebox{\linewidth}{!}{$\displaystyle
\mathit{Aut}(T) = \left\{
\begin{array}{@{}l@{}}
\text{automorphisms } \phi \text{ of } \mu_N \times \Z/N\Z \text{ over } T \text{ which are}\\[2pt]
\text{of the form } {\setlength{\arraycolsep}{2pt}\begin{pmatrix} 1 & * \\ 0 & 1 \end{pmatrix}} \text{ with } * \in \Hom_T(\Z/N\Z, \mu_N)
\end{array}
\right\} .
$}
\]

The assignment $\phi \mapsto *$ defines an isomorphism of $S$-group-schemes
\[
\mathit{Aut} \xrightarrow{\ \sim\ } \mu_N ,
\]
whence it is the group $H^1_{\text{f.p.p.f.}}(S, \mu_N)$ which classifies the $S$-forms of $\mu_N \times \Z/N\Z$ with its extension structure. One verifies easily that given a form $G$, its class in $H^1(S, \mu_N)$ is none other than the isomorphism class of the $\mu_N$-torsor $\pi^{-1}(1)$
\[
0 \to \mu_N \to G \xrightarrow{\ \pi\ } \Z/N\Z \to 0 .
\]

The Kummer sequence on $S$
\[
0 \to \mu_N \to \mathbb{G}_m \xrightarrow{\ N\ } \mathbb{G}_m \to 0
\]
yields a long exact cohomology sequence, the relevant part of which is
\[
\to \mathbb{G}_m(S) \to H^1(S, \mu_N) \to \Pic(S) \to .
\]

Now for any element $q \in \mathbb{G}_m(S)$, the corresponding $\mu_N$-torsor on $S$ is $[N]^{-1}(q)$, i.e., it is the torsor whose $T$-valued points are $\{X \in \mathbb{G}_m(T) \text{ with } X^N = q\}$. Now this is precisely the description of the inverse image of $1$ in
\[
0 \to \mu_N \to \mathbf{T}[N] \to \Z/N\Z \to 0
\]
when we view it as an $S$-group by the extension of scalars $S \to \Spec(\Z[q,q^{-1}])$ which is $q \in \mathbb{G}_m(S)$. Therefore a given $G$ is $S$-isomorphic to a $\mathbf{T}[N]$ for some $q \in \mathbb{G}_m(S)$ if and only if the \emph{image in} $\Pic(S)$ of $\mathrm{cl}(G) \in H^1(S, \mu_N)$ is the zero-element of $\Pic(S)$. But any given element of $\Pic(S)$ dies Zariski locally on $S$.
\end{proof}

\begin{kprop}\label{prop:8.10.6}
Let $G/S$ be a finite locally free commutative group-scheme of rank $N^2$ which is killed by $N$, given with an extension structure
\[
0 \to K \to G \to E \to 0
\]
where $E$ is an etale $S$-group, locally for the etale topology on $S$ isomorphic to $\Z/N\Z$. The following statements (1), (2), (3) are equivalent, and they imply statement (4):
\begin{enumerate}[(1)]
\item The $S$-group-scheme $K$ is \emph{toroidal}, and it is $S$-isomorphic to the \emph{Cartier dual} of $E$.
\item There exists a strongly alternating form $(\ ,\ )$ on $G$ which makes $G$ \emph{auto-dual}.
\item There exists a strongly alternating form $(\ ,\ )$ on $G$ whose induced pairing $\langle\ ,\ \rangle : K \times E \to \mu_N$ makes $K$ and $E$ the \emph{Cartier duals} of each other.
\item Locally for the etale topology on $S$, there exists a unit $q \in \mathbb{G}_m(S)$ and an isomorphism $\mathbf{T}[N] \tens{\Z[q,q^{-1}]} S \xrightarrow{\ \sim\ } G$ respecting the given extension structures.
\end{enumerate}
\end{kprop}

\begin{proof}
Clearly we have $(2) \Longleftrightarrow (3) \Longrightarrow (1)$, and we have $(1) \Longrightarrow (4)$ by the previous proposition. The implication $(1) \Longrightarrow (3)$ holds because already (4) shows that $(\ ,\ )$ exists etale locally on $S$, and it is unique if we require it to induce a given $S$-isomorphism of $K$ with the Cartier dual of $E$. Then by etale descent, we get $(\ ,\ )$ on $G/S$.
\end{proof}

\begin{kcor}\label{cor:8.10.7}
Let $G/S$ be a finite locally free commutative group-scheme of rank $N^2$ which is killed by $N$, which is an extension
\[
0 \to K \to G \to E \to 0
\]
where $E$ is an etale $S$-group, locally (etale) on $S$ isomorphic to $\Z/N\Z$, and where $K$ is given with an $S$-isomorphism to the Cartier dual of $E$. Then
\begin{enumerate}[(1)]
\item there exists a unique strongly alternating form $(\ ,\ )$ which induces, via $\langle\ ,\ \rangle$, the given isomorphism of $K$ with the dual of $E$, and this form $(\ ,\ )$ makes $G$ \emph{auto-dual}.
\item There exists a faithfully flat $S'/S$, an elliptic curve $E'/S'$ and an isomorphism of $S'$-group-schemes
\[
E'[N] \simeq G \underset{S}{\times} S'
\]
which carries the $e_N$ pairing on $E'[N]$ to the pairing $(\ ,\ )$ on $G \underset{S}{\times} S'$.
\end{enumerate}
\end{kcor}

% Section 8.11
\section{Computation of $\widehat{\mathrm{Cusps}}(\mathcal{P})$ via the Tate curve}
\label{sec:8.11}

\kmitem\label{item:8.11.1}
We return to the situation of an excellent noetherian regular ring $R$, and a relatively representable moduli problem $\mathcal{P}$ on $\Ell/R$ which is finite over $\Ell/R$, and for which $\Mcoarse(\mathcal{P})$ is normal near infinity. Then $\overline{\Mcoarse}(\mathcal{P})$ is the unique finite $\PP^1_R$-scheme which is normal near infinity and which agrees with $\Mcoarse(\mathcal{P})$ over $\Aff^1_R$. Consider the cartesian diagram
\begin{equation}\label{eq:8.11.1.1}
\begin{tikzcd}[row sep=1.4em, column sep=0.3em, font=\small]
& & \overline{\Mcoarse}(\mathcal{P}) \arrow[ddd] & & \\
& \widehat{\mathrm{Cusps}}(\mathcal{P}) \arrow[ru, hook] \arrow[ddd] & & \Mcoarse(\mathcal{P}) \arrow[lu, hook'] \arrow[ddd] & \\
\widehat{\mathrm{Cusps}}(\mathcal{P})_{R((1/j))} \arrow[ru, hook] \arrow[ddd] & & & & \Mcoarse(\mathcal{P})_{R((1/j))} \arrow[lu, hook'] \arrow[ddd] \\
& & \PP^1_R & & \\
& \Spec(R[[1/j]]) \arrow[ru, hook] & & \Aff^1_R \arrow[lu, hook'] & \\
\Spec(R((1/j))) \arrow[ru, hook] & & & & \Spec(R((1/j))) \arrow[lu, hook']
\end{tikzcd}
\end{equation}

\begin{klem}\label{lem:8.11.2}
For $\mathcal{P}$ as in \eqref{item:8.11.1} above, $\widehat{\mathrm{Cusps}}(\mathcal{P})$ is the normalization of $R[[1/j]]$ in the normal finite $R((1/j))$-scheme $\Mcoarse(\mathcal{P})_{R((1/j))}$.
\end{klem}

\begin{proof}
Because $R$ is excellent, the morphism
\[
\Spec(R[[1/j]]) \to \Spec(R[1/j])
\]
has all its fibers geometrically regular (cf.~\cite[IV, 7.8.3]{EGA}). Therefore the morphism
\[
\widehat{\mathrm{Cusps}}(\mathcal{P}) \to \overline{\Mcoarse}(\mathcal{P})
\]
has all its fibers geometrically regular. As $\overline{\Mcoarse}(\mathcal{P})$ is normal near infinity, it follows (\cite[VII, 4.8]{AK1}) that $\widehat{\mathrm{Cusps}}(\mathcal{P})$ is normal. As $\widehat{\mathrm{Cusps}}(\mathcal{P})$ is also \emph{finite} over $R[[1/j]]$, it is none other than the \emph{normalization} of $R[[1/j]]$ in $\Mcoarse(\mathcal{P})_{R((1/j))}$.
\end{proof}

\kmitem\label{item:8.11.3}
For any ring $R$, there is a unique $R$-algebra isomorphism
\[
R[[q]] \simeq R[[1/j]]
\]
which is continuous for the adic topologies and which carries
\[
\begin{aligned}
j &\mapsto \frac{1}{q} + 744 + \cdots \\
&= \frac{1}{q} (1 + 744q + \cdots) \\
1/j &\mapsto q(1 - 744q + \cdots) .
\end{aligned}
\]
This isomorphism induces an isomorphism
\[
R((q)) \simeq R((1/j)) .
\]
Expressed in the coordinate $q$, the previous Lemma~\ref{lem:8.11.2} becomes

\begin{klem}\label{lem:8.11.4}
For $\mathcal{P}$ as in \eqref{item:8.11.1} above, the $R[[q]]$-scheme $\widehat{\mathrm{Cusps}}(\mathcal{P})$ is the normalization of $R[[q]]$ in $\Mcoarse(\mathcal{P})_{R((q))}$.
\end{klem}

\kmitem\label{item:8.11.5}
Now view the Tate curve, with its scalars extended to $R((q))$, as an object $\Tate(q)/R((q))$ in $\Ell/R$, and denote by $\mathcal{S}$ the representable moduli problem on $\Ell/R$ which it represents. By the fundamental formula
\[
j(\Tate(q)) = \frac{1}{q} + 744 + \cdots
\]
it is clear that
\[
j(\Tate(q))(j(\Tate(q)) - 1728) \in \frac{1}{q^2}(1 + q R[[q]])
\]
is invertible in $R((q))$.

\kmitem\label{item:8.11.6}
Let $\mathcal{P}$ be any representable moduli problem on $\Ell/R$, and $U$ the open set of $\Aff^1_R$ where $j(j-1728)$ is invertible. We will apply to $\mathcal{P}$ and to $\mathcal{S}$ the result \eqref{cor:8.4.4}. We have just seen that $\mathcal{S} = \mathcal{S}|U$ already lies over $U$, so that
\[
\Mfrak(\mathcal{P}, \mathcal{S}) = \Mfrak(\mathcal{P}, \mathcal{S})|U = \mathcal{P}_{\Tate(q)/R((q))}
\]
\[
(\Mfrak(\mathcal{P}) \underset{\Aff^1_R}{\times} \Mfrak(\mathcal{S}))|U = \Mfrak(\mathcal{P}) \underset{\Aff^1_R}{\times} \Spec(R((q))) ,
\]
and we find

\begin{kprop}\label{prop:8.11.7}
Let $R$ be any ring, $\mathcal{P}$ any representable moduli problem on $\Ell/R$. Then we have a canonical morphism of $R((q))$-schemes
\[
\mathcal{P}_{\Tate(q)/R((q))} \to \Mfrak(\mathcal{P})_{R((q))}
\]
which is a finite etale galois covering with group $\Aut(\Tate(q)/\allowbreak R((q))) = \pm 1$, and so defines an isomorphism of $R((q))$-schemes
\[
(\mathcal{P}_{\Tate(q)/R((q))})/{\pm 1} \xrightarrow{\ \sim\ } \Mfrak(\mathcal{P})_{R((q))} .
\]
\end{kprop}

\begin{kthm}\label{thm:8.11.8}
Let $R$ be any noetherian ring. Then the elliptic curve $\Tate(q)/R((q))$ is \emph{flat} over $\Ell/R$.
\end{kthm}

\begin{proof}
We must show that if $\mathcal{P}$ is any representable moduli problem on $\Ell/R$ which is etale over $\Ell/R$, then the morphism
\[
\mathcal{P}_{\Tate(q)/R((q))} \to \Mfrak(\mathcal{P})
\]
is flat. We may factor this map as
\[
\mathcal{P}_{\Tate(q)/R((q))} \to \Mfrak(\mathcal{P})_{R((q))} \to \Mfrak(\mathcal{P})_{R[j, 1/j]} \to \Mfrak(\mathcal{P}) .
\]
The first map is a finite etale $\pm 1$-torsor, by the preceding proposition. The last two maps are the extension of scalars (by $\Mfrak(\mathcal{P}) \to \Spec(R[j])$) of the maps
\[
\Spec(R((q))) = \Spec(R((1/j))) \to \Spec(R[j, 1/j]) \to \Spec(R[j]) .
\]
The second of these maps is a localization, the first is the extension of scalars (by $R[1/j] \to R[j,1/j]$) of
\[
\Spec(R[[1/j]]) \to \Spec(R[1/j]) ,
\]
which is flat because $R$ is noetherian.
\end{proof}

\begin{kcor}\label{cor:8.11.9}
Let $R$ be a noetherian ring, $\mathcal{P}$ a relatively representable moduli problem on $\Ell/R$ which is affine over $\Ell/R$, and $G$ a finite group operating on $\mathcal{P}$. The canonical morphism
\[
(\mathcal{P}_{\Tate(q)/R((q))})/G \to (\mathcal{P}/G)_{\Tate(q)/R((q))}
\]
is an isomorphism.
\end{kcor}

\begin{proof}
A special case of (\ref{thm:7.1.3}, (3a)), in virtue of the above theorem.
\end{proof}

\begin{kthm}\label{thm:8.11.10}
Let $R$ be an excellent noetherian regular ring, and $\mathcal{P}$ a relatively representable moduli problem on $\Ell/R$ which is finite over $\Ell/R$, and normal near infinity. Suppose either that $\mathcal{P}$ is representable near infinity, or that there exists a prime number $\ell$ which is invertible in $R$. Then the finite $R[[q]]$-scheme $\widehat{\mathrm{Cusps}}(\mathcal{P})$ is the normalization of $R[[q]]$ in the finite normal $R((q))$-scheme
\[
(\mathcal{P}_{\Tate(q)/R((q))})/{\pm 1} .
\]
\end{kthm}

\begin{proof}
If $\mathcal{P}$ is representable near infinity, apply (\ref{prop:8.11.7}) to $\mathcal{P}|U$. If a prime $\ell$ is invertible in $R$, then the moduli problem $[\Gamma(\ell^2)] \tens{\Z} R$ is representable and finite etale galois over $\Ell/R$, with galois group $G = \GL(2, \Z/\ell^2\Z)$. Therefore the moduli problem $\mathcal{P}' = (\mathcal{P}, [\Gamma(\ell^2)] \tens{\Z} R)$ is finite etale galois over $\mathcal{P}$ with the same galois group $G$. Now by the preceding corollary,
\[
(\mathcal{P}'_{\Tate(q)/R((q))})/G \xrightarrow{\ \sim\ } \mathcal{P}_{\Tate(q)/R((q))} ,
\]
whence further dividing by $\pm 1 = \Aut(\Tate(q))$, we find
\[
((\mathcal{P}'_{\Tate(q)/R((q))})/{\pm 1})/G \xrightarrow{\ \sim\ } (\mathcal{P}_{\Tate(q)/R((q))})/{\pm 1} .
\]
Because normalization commutes with passage to quotients by finite groups, we get
\[
\widehat{\mathrm{Cusps}}(\mathcal{P}')/G \xrightarrow{\ \sim\ } \text{normalization of } R[[q]] \text{ in } (\mathcal{P}_{\Tate(q)/R((q))})/{\pm 1}
\]
and the result follows from (\ref{prop:8.6.4}, (3)).
\end{proof}
